
\begin{longtable}[]{@{}
  >{\raggedright\arraybackslash}p{(\linewidth - 2\tabcolsep) * \real{0.5000}}
  >{\raggedright\arraybackslash}p{(\linewidth - 2\tabcolsep) * \real{0.5000}}@{}}
\toprule\noalign{}
\endhead
\bottomrule\noalign{}
\endlastfoot
\multicolumn{2}{@{}>{\raggedright\arraybackslash}p{(\linewidth - 2\tabcolsep) * \real{1.0000} + 2\tabcolsep}@{}}{%
\begin{minipage}[t]{\linewidth}\raggedright
Carl Friedrich Gauss
\end{minipage}} \\
\multicolumn{2}{@{}>{\raggedright\arraybackslash}p{(\linewidth - 2\tabcolsep) * \real{1.0000} + 2\tabcolsep}@{}}{%
\begin{minipage}[t]{\linewidth}\raggedright
{\href{https://en.wikipedia.org/wiki/File:Carl_Friedrich_Gauss_1840_by_Jensen.jpg}{}}

Portrait by \href{https://en.wikipedia.org/wiki/Christian_Albrecht_Jensen}{Christian Albrecht Jensen}, 1840 (copy from Gottlieb Biermann, 1887)
\end{minipage}} \\
Born & \begin{minipage}[t]{\linewidth}\raggedright
Johann Carl Friedrich Gauss

\hfill\break
{({1777-04-30})}30 April 1777\\

\href{https://en.wikipedia.org/wiki/Braunschweig}{Brunswick}, Brunswick-Wolfenbüttel

{}\strut
\end{minipage} \\
Died & \begin{minipage}[t]{\linewidth}\raggedright
23 February 1855{(1855-02-23)} (aged{~}77)\\

\href{https://en.wikipedia.org/wiki/Göttingen}{Göttingen}, Hanover

{}\strut
\end{minipage} \\
Alma{~}mater & \begin{minipage}[t]{\linewidth}\raggedright
\begin{itemize}
\tightlist
\item
  \href{https://en.wikipedia.org/wiki/TU_Braunschweig}{Collegium Carolinum}
\item
  \href{https://en.wikipedia.org/wiki/University_of_Göttingen}{University of Göttingen}
\item
  \href{https://en.wikipedia.org/wiki/University_of_Helmstedt}{University of Helmstedt} (PhD)
\end{itemize}
\end{minipage} \\
Known{~}for & \href{https://en.wikipedia.org/wiki/List_of_things_named_after_Carl_Friedrich_Gauss}{Full list} \\
Spouses & \begin{minipage}[t]{\linewidth}\raggedright
\begin{itemize}
\tightlist
\item
\end{itemize}

Johanna Osthoff

{\hspace{0pt}}

{ }

{\hspace{0pt}}

{(}{m.}{~}1805{;}{ }died{~}1809{)}{\hspace{0pt}}

\begin{itemize}
\tightlist
\item
\end{itemize}

Minna Waldeck

{\hspace{0pt}}

{ }

{\hspace{0pt}}

{(}{m.}{~}1810{;}{ }died{~}1831{)}{\hspace{0pt}}
\end{minipage} \\
Children & 6 \\
Awards & \begin{minipage}[t]{\linewidth}\raggedright
\begin{itemize}
\tightlist
\item
  \href{https://en.wikipedia.org/wiki/Lalande_Prize}{Lalande Prize} (1809)
\item
  \href{https://en.wikipedia.org/wiki/Copley_Medal}{Copley Medal} (1838)
\end{itemize}
\end{minipage} \\
\multicolumn{2}{@{}>{\raggedright\arraybackslash}p{(\linewidth - 2\tabcolsep) * \real{1.0000} + 2\tabcolsep}@{}}{%
\textbf{Scientific career}} \\
Fields & \begin{minipage}[t]{\linewidth}\raggedright
\begin{itemize}
\tightlist
\item
  Mathematics
\item
  astronomy
\item
  geodesy
\item
  magnetism
\end{itemize}
\end{minipage} \\
Institutions & University of Göttingen \\
\href{https://en.wikipedia.org/wiki/Thesis}{Thesis} & \emph{\href{http://www.e-rara.ch/zut/content/titleinfo/1336299}{Demonstratio nova theorematis omnem functionem algebraicam rationalem integram\ldots{}}}{~}{(1799)} \\
\href{https://en.wikipedia.org/wiki/Doctoral_advisor}{Doctoral advisor} & \href{https://en.wikipedia.org/wiki/Johann_Friedrich_Pfaff}{Johann Friedrich Pfaff} \\
Doctoral students & \begin{minipage}[t]{\linewidth}\raggedright
\emph{{~}}

\begin{itemize}
\tightlist
\item
  \href{https://en.wikipedia.org/wiki/Richard_Dedekind}{Richard Dedekind}
\item
  \href{https://en.wikipedia.org/wiki/Christian_Ludwig_Gerling}{Christian Ludwig Gerling}
\item
  \href{https://en.wikipedia.org/wiki/Wilhelm_Klinkerfues}{Wilhelm Klinkerfues}
\item
  \href{https://en.wikipedia.org/wiki/Johann_Benedict_Listing}{Johann Benedict Listing}
\item
  \href{https://en.wikipedia.org/wiki/Bernhard_Riemann}{Bernhard Riemann}
\item
  \href{https://en.wikipedia.org/wiki/August_Ritter_(civil_engineer)}{August Ritter}
\item
  \href{https://en.wikipedia.org/wiki/Karl_von_Staudt}{Karl von Staudt}
\end{itemize}
\end{minipage} \\
Other{~}notable students & \begin{minipage}[t]{\linewidth}\raggedright
\emph{{~}}

\begin{itemize}
\tightlist
\item
  \href{https://en.wikipedia.org/wiki/Gotthold_Eisenstein}{Gotthold Eisenstein}
\item
  \href{https://en.wikipedia.org/wiki/Johann_Franz_Encke}{Johann Franz Encke}
\item
  \href{https://en.wikipedia.org/wiki/Carl_Wolfgang_Benjamin_Goldschmidt}{Carl Wolfgang Benjamin Goldschmidt}
\item
  \href{https://en.wikipedia.org/wiki/Adolph_Theodor_Kupffer}{Adolph Theodor Kupffer}
\item
  \href{https://en.wikipedia.org/wiki/August_Ferdinand_Möbius}{August Ferdinand Möbius}
\item
  \href{https://en.wikipedia.org/wiki/Moritz_Stern}{Moritz Stern}
\item
  \href{https://en.wikipedia.org/wiki/Georg_Frederik_Ursin}{Georg Frederik Ursin}
\item
  \href{https://en.wikipedia.org/wiki/Moritz_Ludwig_George_Wichmann}{Moritz Ludwig Wichmann}
\end{itemize}
\end{minipage} \\
\multicolumn{2}{@{}>{\raggedright\arraybackslash}p{(\linewidth - 2\tabcolsep) * \real{1.0000} + 2\tabcolsep}@{}}{%
} \\
\multicolumn{2}{@{}>{\raggedright\arraybackslash}p{(\linewidth - 2\tabcolsep) * \real{1.0000} + 2\tabcolsep}@{}}{%
Signature} \\
\multicolumn{2}{@{}>{\raggedright\arraybackslash}p{(\linewidth - 2\tabcolsep) * \real{1.0000} + 2\tabcolsep}@{}}{%
{\href{https://en.wikipedia.org/wiki/File:Carl_Friedrich_Gauß_signature.svg}{\includegraphics[width=1.5625in,height=1.39583in]{250px-Carl_Friedrich_Gau_signature.svg.png}}}} \\
\end{longtable}

\textbf{Johann Carl Friedrich Gauss} ({\foreignlanguage{english}{\href{https://en.wikipedia.org/wiki/Help:IPA/English}{/{{ɡ}{aʊ}{s}}/}}{ }{{}\textsuperscript{\href{https://en.wikipedia.org/wiki/File:LL-Q1860_(eng)-Flame,_not_lame-Carl_Friedrich_Gauss.wav}{ⓘ}}}}; {German: }\emph{Gauß}; \foreignlanguage{ngerman}{\href{https://en.wikipedia.org/wiki/Help:IPA/Standard_German}{{{[}}kaʁl{ }ˈfʁiːdʁɪç{ }ˈɡaʊs{{]}}}}{ }{{\href{https://upload.wikimedia.org/wikipedia/commons/transcoded/a/a1/De-carlfriedrichgauss.ogg/De-carlfriedrichgauss.ogg.mp3}{{}{}{}}}\textsuperscript{\href{https://en.wikipedia.org/wiki/File:De-carlfriedrichgauss.ogg}{ⓘ}}}; \href{https://en.wikipedia.org/wiki/Latin_language}{Latin}{: }\emph{Carolus Fridericus Gauss}; 30 April 1777{~}{--}{ }23 February 1855) was a German mathematician, astronomer, geodesist, and physicist, who contributed to many fields in mathematics and science. His mathematical contributions spanned the branches of \href{https://en.wikipedia.org/wiki/Number_theory}{number theory}, \href{https://en.wikipedia.org/wiki/Algebra}{algebra}, \href{https://en.wikipedia.org/wiki/Analysis}{analysis}, \href{https://en.wikipedia.org/wiki/Geometry}{geometry}, \href{https://en.wikipedia.org/wiki/Statistics}{statistics}, and \href{https://en.wikipedia.org/wiki/Probability}{probability}. Gauss was director of the \href{https://en.wikipedia.org/wiki/Göttingen_Observatory}{Göttingen Observatory} in Germany and professor of astronomy from 1807 until his death in 1855.

From an early age, Gauss was known as a \href{https://en.wikipedia.org/wiki/Child_prodigy}{child prodigy} in mathematics. While studying at the \href{https://en.wikipedia.org/wiki/University_of_Göttingen}{University of Göttingen}, he propounded several mathematical \href{https://en.wikipedia.org/wiki/Theorem}{theorems}. As an independent scholar, he wrote the \href{https://en.wikipedia.org/wiki/Masterpiece}{masterpieces} \emph{\href{https://en.wikipedia.org/wiki/Disquisitiones_Arithmeticae}{Disquisitiones Arithmeticae}} and \emph{Theoria motus corporum coelestium}. Gauss produced the second and third complete proofs of the \href{https://en.wikipedia.org/wiki/Fundamental_theorem_of_algebra}{fundamental theorem of algebra}. He also introduced the triple bar symbol (≡) for \href{https://en.wikipedia.org/wiki/Congruence_relation}{congruence}. In \href{https://en.wikipedia.org/wiki/Number_theory}{number theory}, he made numerous contributions, such as the \href{https://en.wikipedia.org/wiki/Gauss_composition_law}{composition law}, the \href{https://en.wikipedia.org/wiki/Quadratic_reciprocity}{law of quadratic reciprocity}, and proved the triangular case of the \href{https://en.wikipedia.org/wiki/Fermat_polygonal_number_theorem}{Fermat polygonal number theorem}. He also contributed to the theory of binary and ternary quadratic forms, and the theory of \href{https://en.wikipedia.org/wiki/Hypergeometric_function}{hypergeometric series}. When Gauss was only 19 years old, he proved the construction of the \href{https://en.wikipedia.org/wiki/Heptadecagon}{heptadecagon}, the first progress in regular polygon construction in over 2000 years. He also introduced the concept of \href{https://en.wikipedia.org/wiki/Gaussian_curvature}{Gaussian curvature} and proved its key properties, especially with his \emph{\href{https://en.wikipedia.org/wiki/Theorema_Egregium}{Theorema Egregium}}. Gauss was the first to prove \href{https://en.wikipedia.org/wiki/Gauss's_inequality}{Gauss\textquotesingle s inequality}. Further, he was instrumental in the development of the \href{https://en.wikipedia.org/wiki/Arithmetic–geometric_mean}{arithmetic--geometric mean}. Due to Gauss\textquotesingle s extensive and fundamental contributions to science and mathematics, more than \href{https://en.wikipedia.org/wiki/List_of_things_named_after_Carl_Friedrich_Gauss}{100 mathematical and scientific concepts} are named after him.

Gauss was instrumental in the identification of \href{https://en.wikipedia.org/wiki/Ceres_(dwarf_planet)}{Ceres} as a dwarf planet. His work on the motion of planetoids disturbed by large planets led to the introduction of the \href{https://en.wikipedia.org/wiki/Gaussian_gravitational_constant}{Gaussian gravitational constant} and the \href{https://en.wikipedia.org/wiki/Least_squares}{method of least squares}, which he had discovered before \href{https://en.wikipedia.org/wiki/Adrien-Marie_Legendre}{Adrien-Marie Legendre} published it. Gauss also introduced the algorithm known as \href{https://en.wikipedia.org/wiki/Recursive_least_squares_filter}{recursive least squares}. Gauss led the geodetic survey of the \href{https://en.wikipedia.org/wiki/Kingdom_of_Hanover}{Kingdom of Hanover} together with an arc measurement project from 1820 to 1844; Gauss was one of the founders of \href{https://en.wikipedia.org/wiki/Geophysics}{geophysics} and formulated the fundamental principles of \href{https://en.wikipedia.org/wiki/Magnetism}{magnetism}. He provided the first absolute measurement of \href{https://en.wikipedia.org/wiki/Earth's_magnetic_field}{Earth\textquotesingle s magnetic field} in 1832, later applying one of his inventions, that of \href{https://en.wikipedia.org/wiki/Spherical_harmonics}{spherical harmonic} analysis, to show that most of Earth\textquotesingle s magnetic field was internal. He was the first to discover and study \href{https://en.wikipedia.org/wiki/Non-Euclidean_geometry}{non-Euclidean geometry}, which he also named. Gauss was the first to develop a \href{https://en.wikipedia.org/wiki/Fast_Fourier_transform}{fast Fourier transform}, doing so some 160 years before \href{https://en.wikipedia.org/wiki/John_Tukey}{John Tukey} and \href{https://en.wikipedia.org/wiki/James_Cooley}{James Cooley}. His practical work led to the invention of the \href{https://en.wikipedia.org/wiki/Heliotrope_(instrument)}{heliotrope} in 1821, a \href{https://en.wikipedia.org/wiki/Magnetometer}{magnetometer} in 1833 and -- with \href{https://en.wikipedia.org/wiki/Wilhelm_Eduard_Weber}{Wilhelm Eduard Weber} -- the first electromagnetic \href{https://en.wikipedia.org/wiki/Telegraphy}{telegraph} in 1833.

Gauss was awarded the \href{https://en.wikipedia.org/wiki/Lalande_Prize}{Lalande Prize} in 1809 for his work on planetary theory and determination of orbits, and the \href{https://en.wikipedia.org/wiki/Copley_Medal}{Copley Medal} in 1838 for his mathematical research in magnetism. He is known for not publishing incomplete work and left several works to be edited \href{https://en.wikipedia.org/wiki/Posthumous_publication}{posthumously}, as a result, this practice delayed the dissemination of many of his discoveries. He believed that the act of learning, not possession of knowledge, provided the greatest enjoyment. While Gauss was not a committed or enthusiastic teacher, generally preferring to focus on his own work, some of his students, such as \href{https://en.wikipedia.org/wiki/Richard_Dedekind}{Richard Dedekind} and \href{https://en.wikipedia.org/wiki/Bernhard_Riemann}{Bernhard Riemann}, became well-known and influential mathematicians in their own right. He married twice and had six children, several of whom later emigrated to the United States.


\subsection{Biography}\label{Biography}


\subsubsection{Youth and education}\label{Youth_and_education}

\begin{figure}
\centering
\href{https://en.wikipedia.org/wiki/File:Braunschweig_Brunswick_Geburtshaus_CF_Gauss_(1914).jpg}{\includegraphics[width=1.97917in,height=2.97917in]{250px-Braunschweig_Brunswick_Geburtshaus_CF_Gauss_1914_.jpg}}
\caption{Birthplace in \href{https://en.wikipedia.org/wiki/Braunschweig}{Brunswick} (destroyed in World War II)}\label{mwcA}
\end{figure}

\begin{figure}
\centering
\href{https://en.wikipedia.org/wiki/File:Goe.Kurze.Geismarstr.Gauss.Wohnhaus.JPG}{\includegraphics[width=1.97917in,height=2.63542in]{250px-Goe.Kurze.Geismarstr.Gauss.Wohnhaus.JPG}}
\caption{Gauss\textquotesingle s home as student in \href{https://en.wikipedia.org/wiki/Göttingen}{Göttingen}}\label{mwdQ}
\end{figure}

Gauss was born on 30 April 1777 in \href{https://en.wikipedia.org/wiki/Braunschweig}{Brunswick}, in the \href{https://en.wikipedia.org/wiki/Duchy_of_Brunswick-Wolfenbüttel}{Duchy of Brunswick-Wolfenbüttel} (now in the German state of \href{https://en.wikipedia.org/wiki/Lower_Saxony}{Lower Saxony}). His family was of relatively low social status. His father Gebhard Dietrich Gauss (1744--1808) worked variously as a butcher, bricklayer, gardener, and treasurer of a death-benefit fund. Gauss characterized his father as honourable and respected, but rough and dominating at home. He was experienced in writing and calculating, whereas his second wife Dorothea, Carl Friedrich\textquotesingle s mother, was nearly illiterate. He had one elder brother from his father\textquotesingle s first marriage.

Gauss was a \href{https://en.wikipedia.org/wiki/Child_prodigy}{child prodigy} in mathematics. When the elementary teachers noticed his intellectual abilities, they brought him to the attention of the \href{https://en.wikipedia.org/wiki/Charles_William_Ferdinand,_Duke_of_Brunswick}{Duke of Brunswick} who sent him to the local \emph{Collegium Carolinum}, which he attended from 1792 to 1795 with \href{https://en.wikipedia.org/wiki/Eberhard_August_Wilhelm_von_Zimmermann}{Eberhard August Wilhelm von Zimmermann} as one of his teachers. Thereafter the Duke granted him the resources for studies of mathematics, sciences, and \href{https://en.wikipedia.org/wiki/Classics}{classical languages} at the \href{https://en.wikipedia.org/wiki/University_of_Göttingen}{University of Göttingen} until 1798. His professor in mathematics was \href{https://en.wikipedia.org/wiki/Abraham_Gotthelf_Kästner}{Abraham Gotthelf Kästner}, whom Gauss called "the leading mathematician among poets, and the leading poet among mathematicians" because of his \href{https://en.wikipedia.org/wiki/Epigram}{epigrams}. Astronomy was taught by \href{https://en.wikipedia.org/wiki/Karl_Felix_Seyffer}{Karl Felix Seyffer}, with whom Gauss stayed in correspondence after graduation; \href{https://en.wikipedia.org/wiki/Heinrich_Wilhelm_Matthias_Olbers}{Olbers} and Gauss mocked him in their correspondence. On the other hand, he thought highly of \href{https://en.wikipedia.org/wiki/Georg_Christoph_Lichtenberg}{Georg Christoph Lichtenberg}, his teacher of physics, and of \href{https://en.wikipedia.org/wiki/Christian_Gottlob_Heyne}{Christian Gottlob Heyne}, whose lectures in classics Gauss attended with pleasure. Fellow students of this time were \href{https://en.wikipedia.org/wiki/Johann_Friedrich_Benzenberg}{Johann Friedrich Benzenberg}, \href{https://en.wikipedia.org/wiki/Farkas_Bolyai}{Farkas Bolyai}, and \href{https://en.wikipedia.org/wiki/Heinrich_Wilhelm_Brandes}{Heinrich Wilhelm Brandes}.

He was likely a self-taught student in mathematics since he independently rediscovered several theorems. He solved a geometrical problem that had occupied mathematicians since the \href{https://en.wikipedia.org/wiki/Ancient_Greece}{Ancient Greeks} when he determined in 1796 which regular \href{https://en.wikipedia.org/wiki/Polygon}{polygons} can be constructed by \href{https://en.wikipedia.org/wiki/Compass_and_straightedge_constructions}{compass and straightedge}. This discovery ultimately led Gauss to choose mathematics instead of \href{https://en.wikipedia.org/wiki/Philology}{philology} as a career. Gauss\textquotesingle s mathematical diary, a collection of short remarks about his results from the years 1796 until 1814, shows that many ideas for his mathematical magnum opus \emph{\href{https://en.wikipedia.org/wiki/Disquisitiones_Arithmeticae}{Disquisitiones Arithmeticae}} (1801) date from this time.

An apocryphal story recounts that as an elementary student, Gauss and his class were tasked by their teacher, J.G. Büttner, to sum the numbers from 1 to 100. Much to Büttner\textquotesingle s surprise, Gauss replied with the correct answer of 5050 in a vastly faster time than expected. Gauss had presumably realised that the sum could be rearranged as 50 pairs of 101 (1 + 100 = 101, 2 + 99= 101, etc.). Thus, he simply multiplied 50 by 101.


\subsubsection{Private scholar}\label{Private_scholar}

Gauss graduated as a \href{https://en.wikipedia.org/wiki/Doctor_of_Philosophy}{Doctor of Philosophy} in 1799, not in Göttingen, as is sometimes stated, but from the \href{https://en.wikipedia.org/wiki/University_of_Helmstedt}{University of Helmstedt}, the only state university of the duchy. \href{https://en.wikipedia.org/wiki/Johann_Friedrich_Pfaff}{Johann Friedrich Pfaff} assessed his doctoral thesis, and Gauss got the degree \emph{\href{https://en.wikipedia.org/wiki/Graduation}{in absentia}} without further oral examination. The Duke then granted him the cost of living as a private scholar in Brunswick. Gauss subsequently refused calls from the \href{https://en.wikipedia.org/wiki/Russian_Academy_of_Sciences}{Russian Academy of Sciences} in \href{https://en.wikipedia.org/wiki/St._Peterburg}{St. Peterburg} and \href{https://en.wikipedia.org/wiki/Ludwig_Maximilian_University_of_Munich}{Landshut University}. Later, the Duke promised him the foundation of an observatory in Brunswick in 1804. Architect \href{https://en.wikipedia.org/wiki/Peter_Joseph_Krahe}{Peter Joseph Krahe} made preliminary designs, but one of \href{https://en.wikipedia.org/wiki/War_of_the_Fourth_Coalition}{Napoleon\textquotesingle s wars} cancelled those plans: the Duke was killed in the \href{https://en.wikipedia.org/wiki/Battle_of_Jena–Auerstedt}{battle of Jena} in 1806. The duchy was abolished in the following year, and Gauss\textquotesingle s financial support stopped.

When Gauss was calculating asteroid orbits in the first years of the century, he established contact with the astronomical communities of \href{https://en.wikipedia.org/wiki/Bremen}{Bremen} and \href{https://en.wikipedia.org/wiki/Lilienthal,_Lower_Saxony}{Lilienthal}, especially \href{https://en.wikipedia.org/wiki/Wilhelm_Olbers}{Wilhelm Olbers}, \href{https://en.wikipedia.org/wiki/Karl_Ludwig_Harding}{Karl Ludwig Harding}, and \href{https://en.wikipedia.org/wiki/Friedrich_Wilhelm_Bessel}{Friedrich Wilhelm Bessel}, forming part of the informal group of astronomers known as the \href{https://en.wikipedia.org/wiki/Celestial_police}{Celestial police}. One of their aims was the discovery of further planets. They assembled data on asteroids and comets as a basis for Gauss\textquotesingle s research on their orbits, which he later published in his astronomical magnum opus \emph{\href{https://en.wikipedia.org/wiki/Theoria_motus_corporum_coelestium?action=edit&redlink=1}{Theoria motus corporum coelestium}} (1809).


\subsubsection[Professor in Göttingen]{\texorpdfstring{\protect\hypertarget{Professor_in_G.C3.B6ttingen}{}{}Professor in Göttingen}{Professor in Göttingen}}\label{Professor_in_Guxf6ttingen}

\begin{figure}
\centering
\href{https://en.wikipedia.org/wiki/File:Goettingen_Sternwarte_01.jpeg}{\includegraphics[width=2.60417in,height=1.70833in]{250px-Goettingen_Sternwarte_01.jpeg}}
\caption{Old Göttingen observatory, {c.}{{ }1800}}\label{mwARk}
\end{figure}

In November 1807, Gauss was hired by the \href{https://en.wikipedia.org/wiki/University_of_Göttingen}{University of Göttingen}, then an institution of the newly founded \href{https://en.wikipedia.org/wiki/Kingdom_of_Westphalia}{Kingdom of Westphalia} under \href{https://en.wikipedia.org/wiki/Jérôme_Bonaparte}{Jérôme Bonaparte}, as full professor and director of the \href{https://en.wikipedia.org/wiki/Göttingen_Observatory}{astronomical observatory}, and kept the chair until his death in 1855. He was soon confronted with the demand for two thousand \href{https://en.wikipedia.org/wiki/Franc}{francs} from the Westphalian government as a war contribution, which he could not afford to pay. Both Olbers and \href{https://en.wikipedia.org/wiki/Pierre-Simon_Laplace}{Laplace} wanted to help him with the payment, but Gauss refused their assistance. Finally, an anonymous person from \href{https://en.wikipedia.org/wiki/Frankfurt}{Frankfurt}, later discovered to be \href{https://en.wikipedia.org/wiki/Prince-primate}{Prince-primate} \href{https://en.wikipedia.org/wiki/Karl_Theodor_Anton_Maria_von_Dalberg}{Dalberg}, paid the sum.

Gauss took on the directorship of the 60-year-old observatory, founded in 1748 by \href{https://en.wikipedia.org/wiki/Prince-elector}{Prince-elector} \href{https://en.wikipedia.org/wiki/George_II_of_Great_Britain}{George II} and built on a converted fortification tower, with usable but partly out-of-date instruments. The construction of a new observatory had been approved by Prince-elector \href{https://en.wikipedia.org/wiki/George_III}{George III} in principle since 1802, and the Westphalian government continued the planning, but Gauss could not move to his new place of work until September 1816. He got new up-to-date instruments, including two \href{https://en.wikipedia.org/wiki/Meridian_circle}{meridian circles} from \href{https://en.wikipedia.org/wiki/Johann_Georg_Repsold}{Repsold} and \href{https://en.wikipedia.org/wiki/Georg_Friedrich_von_Reichenbach}{Reichenbach}, and a \href{https://en.wikipedia.org/wiki/Heliometer}{heliometer} from \href{https://en.wikipedia.org/wiki/Joseph_von_Fraunhofer}{Fraunhofer}.

The scientific activity of Gauss, besides pure mathematics, can be roughly divided into three periods: astronomy was the main focus in the first two decades of the 19th century, geodesy in the third decade, and physics, mainly magnetism, in the fourth decade.

Gauss made no secret of his aversion to giving academic lectures. But from the start of his academic career at Göttingen, he continuously gave lectures until 1854. He often complained about the burdens of teaching, feeling that it was a waste of his time. On the other hand, he occasionally described some students as talented. Most of his lectures dealt with astronomy, geodesy, and \href{https://en.wikipedia.org/wiki/Applied_mathematics}{applied mathematics}, and only three lectures on subjects of pure mathematics. Some of Gauss\textquotesingle s students went on to become renowned mathematicians, physicists, and astronomers: \href{https://en.wikipedia.org/wiki/Moritz_Cantor}{Moritz Cantor}, \href{https://en.wikipedia.org/wiki/Richard_Dedekind}{Dedekind}, \href{https://en.wikipedia.org/wiki/Enno_Dirksen}{Dirksen}, \href{https://en.wikipedia.org/wiki/Johann_Franz_Encke}{Encke}, \href{https://en.wikipedia.org/wiki/Benjamin_Apthorp_Gould}{Gould}, \href{https://en.wikipedia.org/wiki/Eduard_Heine}{Heine}, \href{https://en.wikipedia.org/wiki/Ernst_Klinkerfues}{Klinkerfues}, \href{https://en.wikipedia.org/wiki/Adolph_Theodor_Kupffer}{Kupffer}, \href{https://en.wikipedia.org/wiki/Johann_Benedict_Listing}{Listing}, \href{https://en.wikipedia.org/wiki/August_Ferdinand_Möbius}{Möbius}, \href{https://en.wikipedia.org/wiki/Friedrich_Bernhard_Gottfried_Nicolai}{Nicolai}, \href{https://en.wikipedia.org/wiki/Bernhard_Riemann}{Riemann}, \href{https://en.wikipedia.org/wiki/August_Ritter_(civil_engineer)}{Ritter}, \href{https://en.wikipedia.org/wiki/Ernst_Christian_Julius_Schering}{Schering}, \href{https://en.wikipedia.org/wiki/Heinrich_Scherk}{Scherk}, \href{https://en.wikipedia.org/wiki/Heinrich_Christian_Schumacher}{Schumacher}, \href{https://en.wikipedia.org/wiki/Karl_Georg_Christian_von_Staudt}{von Staudt}, \href{https://en.wikipedia.org/wiki/Moritz_Abraham_Stern}{Stern}, \href{https://en.wikipedia.org/wiki/Georg_Frederik_Ursin}{Ursin}; as geoscientists \href{https://en.wikipedia.org/wiki/Wolfgang_Sartorius_von_Waltershausen}{Sartorius von Waltershausen}, and \href{https://en.wikipedia.org/wiki/Johann_Eduard_Wappäus}{Wappäus}.

Gauss did not write any textbook and disliked the \href{https://en.wikipedia.org/wiki/Popular_science}{popularization} of scientific matters. His only attempts at popularization were his works on the date of Easter (1800/1802) and the essay \emph{Erdmagnetismus und Magnetometer} of 1836. Gauss published his papers and books exclusively in \href{https://en.wikipedia.org/wiki/Latin}{Latin} or \href{https://en.wikipedia.org/wiki/German_language}{German}. He wrote Latin in a classical style but used some customary modifications set by contemporary mathematicians.

\begin{figure}
\centering
\href{https://en.wikipedia.org/wiki/File:Universitäts-Sternwarte_Göttingen_02.jpg}{\includegraphics[width=2.60417in,height=1.92708in]{250px-Universit_ts-Sternwarte_G_ttingen_02.jpg}}
\caption{The new Göttingen observatory of 1816; Gauss\textquotesingle s living rooms were in the western wing (right)}\label{mwAa8}
\end{figure}

\begin{figure}
\centering
\href{https://en.wikipedia.org/wiki/File:Die_Göttinger_Sieben_von_Eduard_Ritmüller.jpg}{\includegraphics[width=1.97917in,height=2.84375in]{250px-Die_G_ttinger_Sieben_von_Eduard_Ritm_ller.jpg}}
\caption{\href{https://en.wikipedia.org/wiki/Wilhelm_Eduard_Weber}{Wilhelm Weber} and \href{https://en.wikipedia.org/wiki/Heinrich_Ewald}{Heinrich Ewald} (first row) as members of the \href{https://en.wikipedia.org/wiki/Göttingen_Seven}{Göttingen Seven}}\label{mwAbM}
\end{figure}

\begin{figure}
\centering
\href{https://en.wikipedia.org/wiki/File:Carl_Friedrich_Gauss_on_his_Deathbed,_1855.jpg}{\includegraphics[width=2.60417in,height=1.92708in]{250px-Carl_Friedrich_Gauss_on_his_Deathbed_1855.jpg}}
\caption{Gauss on his deathbed (1855) (daguerreotype from Philipp Petri)}\label{mwAbo}
\end{figure}

Gauss gave his inaugural lecture at Göttingen University in 1808. He described his approach to astronomy as based on reliable observations and accurate calculations, rather than on belief or empty hypothesizing. At university, he was accompanied by a staff of other lecturers in his disciplines, who completed the educational program; these included the mathematician Thibaut with his lectures, the physicist \href{https://en.wikipedia.org/wiki/Johann_Tobias_Mayer}{Mayer}, known for his textbooks, his successor \href{https://en.wikipedia.org/wiki/Wilhelm_Eduard_Weber}{Weber} since 1831, and in the observatory \href{https://en.wikipedia.org/wiki/Karl_Ludwig_Harding}{Harding}, who took the main part of lectures in practical astronomy. When the observatory was completed, Gauss occupied the western wing of the new observatory, while Harding took the eastern. They had once been on friendly terms, but over time they became alienated, possibly -- as some biographers presume -- because Gauss had wished the equal-ranked Harding to be no more than his assistant or observer. Gauss used the new \href{https://en.wikipedia.org/wiki/Meridian_circle}{meridian circles} nearly exclusively, and kept them away from Harding, except for some very seldom joint observations.

\href{https://en.wikipedia.org/wiki/Martin_Brendel}{Brendel} subdivides Gauss\textquotesingle s astronomic activity chronologically into seven periods, of which the years since 1820 are taken as a "period of lower astronomical activity". The new, well-equipped observatory did not work as effectively as other ones; Gauss\textquotesingle s astronomical research had the character of a one-man enterprise without a long-time observation program, and the university established a place for an assistant only after Harding died in 1834.

Nevertheless, Gauss twice refused the opportunity to solve the problem, turning down offers from Berlin in 1810 and 1825 to become a full member of the Prussian Academy without burdening lecturing duties, as well as from \href{https://en.wikipedia.org/wiki/Leipzig_University}{Leipzig University} in 1810 and from \href{https://en.wikipedia.org/wiki/Vienna_University}{Vienna University} in 1842, perhaps because of the family\textquotesingle s difficult situation. Gauss\textquotesingle s salary was raised from 1000 \href{https://en.wikipedia.org/wiki/Reichsthaler}{Reichsthaler} in 1810 to 2500 Reichsthaler in 1824, and in his later years he was one of the best-paid professors of the university.

When Gauss was asked for help by his colleague and friend \href{https://en.wikipedia.org/wiki/Friedrich_Wilhelm_Bessel}{Friedrich Wilhelm Bessel} in 1810, who was in trouble at \href{https://en.wikipedia.org/wiki/Königsberg_University}{Königsberg University} because of his lack of an academic title, Gauss provided a \href{https://en.wikipedia.org/wiki/Honorary_degree}{doctorate \emph{honoris causa}} for Bessel from the Philosophy Faculty of Göttingen in March 1811. Gauss gave another recommendation for an honorary degree for \href{https://en.wikipedia.org/wiki/Sophie_Germain}{Sophie Germain} but only shortly before her death, so she never received it. He also gave successful support to the mathematician \href{https://en.wikipedia.org/wiki/Gotthold_Eisenstein}{Gotthold Eisenstein} in Berlin.

Gauss was loyal to the \href{https://en.wikipedia.org/wiki/House_of_Hanover}{House of Hanover}. After King \href{https://en.wikipedia.org/wiki/William_IV}{William IV} died in 1837, the new Hanoverian King \href{https://en.wikipedia.org/wiki/Ernest_Augustus,_King_of_Hanover}{Ernest Augustus} annulled the 1833 constitution. Seven professors, later known as the "\href{https://en.wikipedia.org/wiki/Göttingen_Seven}{Göttingen Seven}", protested against this, among them his friend and collaborator Wilhelm Weber and Gauss\textquotesingle s son-in-law Heinrich Ewald. All of them were dismissed, and three of them were expelled, but Ewald and Weber could stay in Göttingen. Gauss was deeply affected by this quarrel but saw no possibility to help them.

Gauss took part in academic administration: three times he was elected as \href{https://en.wikipedia.org/wiki/Dean_(education)}{dean} of the Faculty of Philosophy. Being entrusted with the widow\textquotesingle s \href{https://en.wikipedia.org/wiki/Pension_fund}{pension fund} of the university, he dealt with \href{https://en.wikipedia.org/wiki/Actuarial_science}{actuarial science} and wrote a report on the strategy for stabilizing the benefits. He was appointed director of the Royal Academy of Sciences in Göttingen for nine years.

Gauss remained mentally active into his old age, even while suffering from \href{https://en.wikipedia.org/wiki/Gout}{gout} and general unhappiness. On 23 February 1855, he died of a heart attack in Göttingen; and was interred in the \href{https://en.wikipedia.org/wiki/Albanifriedhof}{Albani Cemetery} there. \href{https://en.wikipedia.org/wiki/Heinrich_Ewald}{Heinrich Ewald}, Gauss\textquotesingle s son-in-law, and \href{https://en.wikipedia.org/wiki/Wolfgang_Sartorius_von_Waltershausen}{Wolfgang Sartorius von Waltershausen}, Gauss\textquotesingle s close friend and biographer, gave eulogies at his funeral.

Gauss was a successful investor and accumulated considerable wealth with stocks and securities, amounting to a value of more than 150,000 Thaler; after his death, about 18,000 Thaler were found hidden in his rooms.


\subsubsection[Gauss\textquotesingle s brain]{\texorpdfstring{\protect\hypertarget{Gauss.27s_brain}{}{}Gauss\textquotesingle s brain}{Gauss\textquotesingle s brain}}\label{Gaussux27s_brain}

The day after Gauss\textquotesingle s death his brain was removed, preserved, and studied by \href{https://en.wikipedia.org/wiki/Rudolf_Wagner}{Rudolf Wagner}, who found its mass to be slightly above average, at {1,492 grams (3.29}{~}{lb)}. Wagner\textquotesingle s son \href{https://en.wikipedia.org/wiki/Hermann_Wagner_(geographer)}{Hermann}, a geographer, estimated the cerebral area to be {219,588 square millimetres (340.362}{~}{sq}{~}{in)} in his doctoral thesis. In 2013, a neurobiologist at the \href{https://en.wikipedia.org/wiki/Max_Planck_Institute_for_Biophysical_Chemistry}{Max Planck Institute for Biophysical Chemistry} in Göttingen discovered that Gauss\textquotesingle s brain had been mixed up soon after the first investigations, due to mislabelling, with that of the physician \href{https://en.wikipedia.org/wiki/Conrad_Heinrich_Fuchs}{Conrad Heinrich Fuchs}, who died in Göttingen a few months after Gauss. A further investigation showed no remarkable anomalies in the brains of either person. Thus, all investigations of Gauss\textquotesingle s brain until 1998, except the first ones of Rudolf and Hermann Wagner, actually refer to the brain of Fuchs.


\subsubsection{Family}\label{Family}

\begin{figure}
\centering
\href{https://en.wikipedia.org/wiki/File:Minna_Gauß_geb._Waldeck,_002.jpg}{\includegraphics[width=1.97917in,height=2.88542in]{250px-Minna_Gau_geb._Waldeck_002.jpg}}
\caption{Gauss\textquotesingle s second wife Wilhelmine Waldeck}\label{mwAlw}
\end{figure}

Gauss married Johanna Osthoff on 9 October 1805 in St. Catherine\textquotesingle s church in Brunswick. They had two sons and one daughter: Joseph (1806--1873), Wilhelmina (1808--1840), and Louis (1809--1810). Johanna died on 11 October 1809, one month after the birth of Louis, who himself died a few months later. Gauss chose the first names of his children in honour of \href{https://en.wikipedia.org/wiki/Giuseppe_Piazzi}{Giuseppe Piazzi}, Wilhelm Olbers, and Karl Ludwig Harding, the discoverers of the first asteroids.

On 4 August 1810, Gauss married Wilhelmine (Minna) Waldeck, a friend of his first wife, with whom he had three more children: Eugen (later Eugene) (1811--1896), Wilhelm (later William) (1813--1879), and Therese (1816--1864). Minna Gauss died on 12 September 1831 after being seriously ill for more than a decade. Therese then took over the household and cared for Gauss for the rest of his life; after her father\textquotesingle s death, she married actor Constantin Staufenau. Her sister Wilhelmina married the orientalist \href{https://en.wikipedia.org/wiki/Heinrich_Ewald}{Heinrich Ewald}. Gauss\textquotesingle s mother Dorothea lived in his house from 1817 until she died in 1839.

The eldest son Joseph, while still a schoolboy, helped his father as an assistant during the survey campaign in the summer of 1821. After a short time at university, in 1824 Joseph joined the \href{https://en.wikipedia.org/wiki/Hanoverian_army}{Hanoverian army} and assisted in surveying again in 1829. In the 1830s he was responsible for the enlargement of the survey network into the western parts of the kingdom. With his geodetical qualifications, he left the service and engaged in the construction of the railway network as director of the \href{https://en.wikipedia.org/wiki/Royal_Hanoverian_State_Railways}{Royal Hanoverian State Railways}. In 1836 he studied the railroad system in the US for some months.

Eugen left Göttingen in September 1830 and emigrated to the United States, where he spent five years with the army. He then worked for the \href{https://en.wikipedia.org/wiki/American_Fur_Company}{American Fur Company} in the Midwest. He later moved to \href{https://en.wikipedia.org/wiki/Missouri}{Missouri} and became a successful businessman. Wilhelm married a niece of the astronomer \href{https://en.wikipedia.org/wiki/Friedrich_Bessel}{Bessel}; he then moved to Missouri, started as a farmer and became wealthy in the shoe business in \href{https://en.wikipedia.org/wiki/St._Louis,_Missouri}{St. Louis} in later years. Eugene and William have numerous descendants in America, but the Gauss descendants left in Germany all derive from Joseph, as the daughters had no children.

\begin{itemize}
\item
  
  
  {\href{https://en.wikipedia.org/wiki/File:Joseph_Gauß,_001.jpg}{\includegraphics[width=0.67708in,height=1.25in]{120px-Joseph_Gau_001.jpg}}}

  
  Joseph Gauss
\item
  
  
  {\href{https://en.wikipedia.org/wiki/File:Joseph_Gauß,_003.jpg}{\includegraphics[width=0.97917in,height=1.25in]{120px-Joseph_Gau_003.jpg}}}

  
  Sophie Gauss née Erythropel\\
  Joseph\textquotesingle s wife
\item
  
  
  {\href{https://en.wikipedia.org/wiki/File:Minna_Ewald_geb._Gauß,_003.jpg}{\includegraphics[width=1.02083in,height=1.25in]{120px-Minna_Ewald_geb._Gau_003.jpg}}}

  
  Wilhelmina Gauss
\item
  
  
  {\href{https://en.wikipedia.org/wiki/File:Ewald,_Georg_Heinrich_August_(1803-1875).jpg}{\includegraphics[width=0.86458in,height=1.25in]{120px-Ewald_Georg_Heinrich_August_1803-1875_.jpg}}}

  
  \href{https://en.wikipedia.org/wiki/Heinrich_Ewald}{Heinrich Ewald}\\
  Wilhelmina\textquotesingle s husband
\item
  
  
  {\href{https://en.wikipedia.org/wiki/File:Eugen_Gauß,_001.jpg}{\includegraphics[width=0.8125in,height=1.25in]{120px-Eugen_Gau_001.jpg}}}

  
  Eugen (Eugene) Gauss
\item
  
  
  {\href{https://en.wikipedia.org/wiki/File:Eugen_Gauß,_003.jpg}{\includegraphics[width=0.875in,height=1.25in]{120px-Eugen_Gau_003.jpg}}}

  
  Henrietta Gauss née Fawcett\\
  Eugene\textquotesingle s wife
\item
  
  
  {\href{https://en.wikipedia.org/wiki/File:Wilhelm_Gauß,_002.jpg}{\includegraphics[width=0.79167in,height=1.25in]{120px-Wilhelm_Gau_002.jpg}}}

  
  Wilhelm (Charles William) Gauss
\item
  
  
  {\href{https://en.wikipedia.org/wiki/File:Wilhelm_Gauß,_001.jpg}{\includegraphics[width=0.9375in,height=1.25in]{120px-Wilhelm_Gau_001.jpg}}}

  
  Louisa Aletta Gauss née Fallenstein\\
  William\textquotesingle s wife
\item
  
  
  {\href{https://en.wikipedia.org/wiki/File:Therese_Staufenau_geb._Gauß,_008.jpg}{\includegraphics[width=1.04167in,height=1.25in]{120px-Therese_Staufenau_geb._Gau_008.jpg}}}

  
  Therese Gauss
\item
  
  
  {\href{https://en.wikipedia.org/wiki/File:Therese_Staufenau_geb._Gauß,_010.jpg}{\includegraphics[width=0.83333in,height=1.25in]{120px-Therese_Staufenau_geb._Gau_010.jpg}}}

  
  Constantin Staufenau\\
  Therese\textquotesingle s husband
\end{itemize}


\subsubsection{Personality}\label{Personality}


\paragraph{Scholar}\label{Scholar}

\begin{figure}
\centering
\href{https://en.wikipedia.org/wiki/File:Carl_Friedrich_Gauß,_Karikatur_von_Abraham_Gotthelf_Kästner,_1795.jpg}{\includegraphics[width=2.60417in,height=1.90625in]{250px-Carl_Friedrich_Gau_Karikatur_von_Abraham_Gotthelf_K_stner_1795.jpg}}
\caption{A student draws his professor of mathematics: Caricature of \href{https://en.wikipedia.org/wiki/Abraham_Gotthelf_Kästner}{Abraham Gotthelf Kästner} by Gauss (1795)}\label{mwAuU}
\end{figure}

\begin{figure}
\centering
\href{https://en.wikipedia.org/wiki/File:Carl_Friedrich_Gauß,_003.jpg}{\includegraphics[width=1.97917in,height=2.79167in]{250px-Carl_Friedrich_Gau_003.jpg}}
\caption{A student draws his professor of mathematics: Gauss sketched by his student \href{https://en.wikipedia.org/wiki/Johann_Benedict_Listing}{Johann Benedict Listing}, 1830}\label{mwAu4}
\end{figure}

In the first two decades of the 19th century, Gauss was the only important mathematician in Germany comparable to the leading French mathematicians. His \emph{Disquisitiones Arithmeticae} was the first mathematical book from Germany to be translated into the French language.

Gauss was "in front of the new development" with documented research since 1799, his wealth of new ideas, and his rigour of demonstration. In contrast to previous mathematicians like \href{https://en.wikipedia.org/wiki/Leonhard_Euler}{Leonhard Euler}, who let their readers take part in their reasoning, including certain erroneous deviations from the correct path, Gauss introduced a new style of direct and complete exposition that did not attempt to show the reader the author\textquotesingle s train of thought.

\begin{quote}
Gauss was the first to restore that \emph{rigor} of demonstration which we admire in the ancients and which had been forced unduly into the background by the exclusive interest of the preceding period in \emph{new} developments.
\end{quote}


---{ }\href{https://en.wikipedia.org/wiki/Carl_Friedrich_Gauss\#CITEREFKlein1894}{Klein 1894}, p.{~}101{}

But for himself, he propagated a quite different ideal, given in a letter to Farkas Bolyai as follows:

\begin{quote}
It is not knowledge, but the act of learning, not possession but the act of getting there, which grants the greatest enjoyment. When I have clarified and exhausted a subject, then I turn away from it, in order to go into darkness again.
\end{quote}


---{ }\href{https://en.wikipedia.org/wiki/Carl_Friedrich_Gauss\#CITEREFDunnington2004}{Dunnington 2004}, p.{~}416{}

His posthumous papers, his scientific \href{https://en.wikipedia.org/wiki/Gauss's_diary}{diary}, and short glosses in his own textbooks show that he empirically worked to a great extent. He was a lifelong busy and enthusiastic calculator, working extraordinarily quickly and checking his results through estimation. Nevertheless, his calculations were not always free from mistakes. He coped with the enormous workload by using skillful tools. Gauss used numerous \href{https://en.wikipedia.org/wiki/Mathematical_table}{mathematical tables}, examined their exactness, and constructed new tables on various matters for personal use. He developed new tools for effective calculation, for example the \href{https://en.wikipedia.org/wiki/Gaussian_elimination}{Gaussian elimination}. Gauss\textquotesingle s calculations and the tables he prepared were often more precise than practically necessary. Very likely, this method gave him additional material for his theoretical work.

\begin{figure}
\centering
\href{https://en.wikipedia.org/wiki/File:GaussSiegel1777-1855.png}{
\includegraphics[width=1.97917in,height=1.97917in]{250px-GaussSiegel1777-1855.png}}
\caption{Gauss\textquotesingle s \href{https://en.wikipedia.org/wiki/Seal_(emblem)}{seal} with his motto \emph{Pauca sed Matura}}\label{mwA0A}
\end{figure}

Gauss was only willing to publish work when he considered it complete and above criticism. This \href{https://en.wikipedia.org/wiki/Perfectionism_(psychology)}{perfectionism} was in keeping with the motto of his personal \href{https://en.wikipedia.org/wiki/Seal_(emblem)}{seal} {\emph{Pauca sed Matura}} ("Few, but Ripe"). Many colleagues encouraged him to publicize new ideas and sometimes rebuked him if he hesitated too long, in their opinion. Gauss defended himself by claiming that the initial discovery of ideas was easy, but preparing a presentable elaboration was a demanding matter for him, for either lack of time or "serenity of mind". Nevertheless, he published many short communications of urgent content in various journals, but left a considerable literary estate, too. Gauss referred to mathematics as "the queen of sciences" and arithmetics as "the queen of mathematics", and supposedly once espoused a belief in the necessity of immediately understanding \href{https://en.wikipedia.org/wiki/Euler's_identity}{Euler\textquotesingle s identity} as a benchmark pursuant to becoming a first-class mathematician.

On certain occasions, Gauss claimed that the ideas of another scholar had already been in his possession previously. Thus his concept of priority as "the first to discover, not the first to publish" differed from that of his scientific contemporaries. In contrast to his perfectionism in presenting mathematical ideas, his citations were criticized as negligent. He justified himself with an unusual view of correct citation practice: he would only give complete references, with respect to the previous authors of importance, which no one should ignore, but citing in this way would require knowledge of the history of science and more time than he wished to spend.


\paragraph{Private man}\label{Private_man}

Soon after Gauss\textquotesingle s death, his friend Sartorius published the first biography (1856), written in a rather enthusiastic style. Sartorius saw him as a serene and forward-striving man with childlike modesty, but also of "iron character" with an unshakeable strength of mind. Apart from his closer circle, others regarded him as reserved and unapproachable "like an \href{https://en.wikipedia.org/wiki/Twelve_Olympians}{Olympian} sitting enthroned on the summit of science". His close contemporaries agreed that Gauss was a man of difficult character. He often refused to accept compliments. His visitors were occasionally irritated by his grumpy behaviour, but a short time later his mood could change, and he would become a charming, open-minded host. Gauss disliked polemic natures; together with his colleague \href{https://en.wikipedia.org/wiki/Johann_Friedrich_Ludwig_Hausmann}{Hausmann} he opposed to a call for \href{https://en.wikipedia.org/wiki/Justus_Liebig}{Justus Liebig} on a university chair in Göttingen, "because he was always involved in some polemic."

\begin{figure}
\centering
\href{https://en.wikipedia.org/wiki/File:Göttingen,_Kurze_Straße_15,_001.jpg}{\includegraphics[width=2.60417in,height=1.95833in]{250px-G_ttingen_Kurze_Stra_e_15_001.jpg}}
\caption{Gauss\textquotesingle s residence from 1808 to 1816 in the first floor}\label{mwA4Y}
\end{figure}

Gauss\textquotesingle s life was overshadowed by severe problems in his family. When his first wife Johanna suddenly died shortly after the birth of their third child, he revealed the grief in a last letter to his dead wife in the style of an ancient \href{https://en.wikipedia.org/wiki/Threnody}{threnody}, the most personal of his surviving documents. His second wife and his two daughters suffered from \href{https://en.wikipedia.org/wiki/Tuberculosis}{tuberculosis}. In a letter to \href{https://en.wikipedia.org/wiki/Friedrich_Wilhelm_Bessel}{Bessel}, dated December 1831, Gauss hinted at his distress, describing himself as "the victim of the worst domestic sufferings".

Because of his wife\textquotesingle s illness, both younger sons were educated for some years in \href{https://en.wikipedia.org/wiki/Celle}{Celle}, far from Göttingen. The military career of his elder son Joseph ended after more than two decades at the poorly paid rank of \href{https://en.wikipedia.org/wiki/First_lieutenant}{first lieutenant}, although he had acquired a considerable knowledge of geodesy. He needed financial support from his father even after he was married. The second son Eugen shared a good measure of his father\textquotesingle s talent in computation and languages but had a lively and sometimes rebellious character. He wanted to study philology, whereas Gauss wanted him to become a lawyer. Having run up debts and caused a scandal in public, Eugen suddenly left Göttingen under dramatic circumstances in September 1830 and emigrated via Bremen to the United States. He wasted the little money he had taken to start, after which his father refused further financial support. The youngest son Wilhelm wanted to qualify for agricultural administration, but had difficulties getting an appropriate education, and eventually emigrated as well. Only Gauss\textquotesingle s youngest daughter Therese accompanied him in his last years of life.

In his later years Gauss habitually collected various types of useful or useless numerical data, such as the number of paths from his home to certain places in Göttingen or peoples\textquotesingle{} ages in days; he congratulated \href{https://en.wikipedia.org/wiki/Alexander_von_Humboldt}{Humboldt} in December 1851 for having reached the same age as \href{https://en.wikipedia.org/wiki/Isaac_Newton}{Isaac Newton} at his death, calculated in days.

Beyond his excellent knowledge of \href{https://en.wikipedia.org/wiki/Latin}{Latin}, he was also acquainted with modern languages. Gauss read both classical and modern literature, and English and French works in the original languages. His favorite English author was \href{https://en.wikipedia.org/wiki/Walter_Scott}{Walter Scott}, his favorite German \href{https://en.wikipedia.org/wiki/Jean_Paul}{Jean Paul}. At the age of 62, he began to teach himself \href{https://en.wikipedia.org/wiki/Russian_language}{Russian}, very likely to understand scientific writings from Russia, among them those of \href{https://en.wikipedia.org/wiki/Nikolai_Lobachevsky}{Lobachevsky} on non-Euclidean geometry. Gauss liked singing and went to concerts. He was a busy newspaper reader; in his last years, he would visit an academic press salon of the university every noon. Gauss did not care much for philosophy, and mocked the "splitting hairs of the so-called metaphysicians", by which he meant proponents of the contemporary school of \emph{\href{https://en.wikipedia.org/wiki/Naturphilosophie}{Naturphilosophie}}.

Gauss had an "aristocratic and through and through conservative nature", with little respect for people\textquotesingle s intelligence and morals, following the motto "\href{https://en.wikipedia.org/wiki/Mundus_vult_decipi,_ergo_decipiatur}{mundus vult decipi}". He disliked Napoleon and his system and was horrified by violence and revolution of all kinds. Thus he condemned the methods of the \href{https://en.wikipedia.org/wiki/Revolutions_of_1848}{Revolutions of 1848}, though he agreed with some of their aims, such as that of a unified Germany. He had a low estimation of the constitutional system and he criticized parliamentarians of his time for their perceived ignorance and logical errors.

Some Gauss biographers have speculated on his religious beliefs. He sometimes said "God arithmetizes" and "I succeeded -- not on account of my hard efforts, but by the grace of the Lord." Gauss was a member of the \href{https://en.wikipedia.org/wiki/Lutheran_church}{Lutheran church}, like most of the population in northern Germany, but it seems that he did not believe all Lutheran \href{https://en.wikipedia.org/wiki/Dogma}{dogma} or understand the Bible fully literally. According to Sartorius, Gauss\textquotesingle s \href{https://en.wikipedia.org/wiki/Religious_tolerance}{religious tolerance}, "insatiable thirst for truth" and sense of justice were motivated by his religious convictions.


\subsection{Mathematics}\label{Mathematics}


\subsubsection{Algebra and number theory}\label{Algebra_and_number_theory}


\paragraph{Fundamental theorem of algebra}\label{Fundamental_theorem_of_algebra}

\begin{figure}
\centering
\href{https://en.wikipedia.org/wiki/File:DBP_1977_928_Carl_Friedrich_Gauß.jpg}{\includegraphics[width=1.97917in,height=3.15625in]{250px-DBP_1977_928_Carl_Friedrich_Gau_.jpg}}
\caption{German stamp commemorating Gauss\textquotesingle s 200th anniversary: the \href{https://en.wikipedia.org/wiki/Complex_plane}{complex plane} or \emph{Gauss plane}}\label{mwBAY}
\end{figure}

In his doctoral thesis from 1799, Gauss proved the \href{https://en.wikipedia.org/wiki/Fundamental_theorem_of_algebra}{fundamental theorem of algebra} which states that every non-constant single-variable \href{https://en.wikipedia.org/wiki/Polynomial}{polynomial} with complex coefficients has at least one complex \href{https://en.wikipedia.org/wiki/Root_of_a_function}{root}. Mathematicians including \href{https://en.wikipedia.org/wiki/Jean_le_Rond_d'Alembert}{Jean le Rond d\textquotesingle Alembert} had produced false proofs before him, and Gauss\textquotesingle s dissertation contains a critique of d\textquotesingle Alembert\textquotesingle s work. He subsequently produced three other proofs, the last one in 1849 being generally rigorous. His attempts led to considerable clarification of the concept of complex numbers.


\paragraph{\texorpdfstring{\emph{Disquisitiones Arithmeticae}}{Disquisitiones Arithmeticae}}\label{Disquisitiones_Arithmeticae}

In the preface to the \emph{Disquisitiones}, Gauss dates the beginning of his work on number theory to 1795. By studying the works of previous mathematicians like \href{https://en.wikipedia.org/wiki/Pierre_de_Fermat}{Fermat}, Euler, Lagrange, and Legendre, he realized that these scholars had already found much of what he had independently discovered. The \emph{\href{https://en.wikipedia.org/wiki/Disquisitiones_Arithmeticae}{Disquisitiones Arithmeticae}}, written in 1798 and published in 1801, consolidated number theory as a discipline and covered both elementary and algebraic \href{https://en.wikipedia.org/wiki/Number_theory}{number theory}. Therein he introduces the \href{https://en.wikipedia.org/wiki/Triple_bar}{triple bar} symbol ({≡}) for \href{https://en.wikipedia.org/wiki/Congruence_relation}{congruence} and uses it for a clean presentation of \href{https://en.wikipedia.org/wiki/Modular_arithmetic}{modular arithmetic}. It deals with the \href{https://en.wikipedia.org/wiki/Unique_factorization_theorem}{unique factorization theorem} and \href{https://en.wikipedia.org/wiki/Primitive_root_modulo_n}{primitive roots modulo n}. In the main sections, Gauss presents the first two proofs of the law of \href{https://en.wikipedia.org/wiki/Quadratic_reciprocity}{quadratic reciprocity} and develops the theories of \href{https://en.wikipedia.org/wiki/Binary_quadratic_form}{binary} and ternary \href{https://en.wikipedia.org/wiki/Quadratic_form}{quadratic forms}.

The \emph{Disquisitiones} include the \href{https://en.wikipedia.org/wiki/Gauss_composition_law}{Gauss composition law} for binary quadratic forms, as well as the enumeration of the number of representations of an integer as the sum of three squares. As an almost immediate corollary of his \href{https://en.wikipedia.org/wiki/Legendre's_three-square_theorem}{theorem on three squares}, he proves the triangular case of the \href{https://en.wikipedia.org/wiki/Fermat_polygonal_number_theorem}{Fermat polygonal number theorem} for \emph{n} = 3. From several analytic results on \href{https://en.wikipedia.org/wiki/Ideal_class_group}{class numbers} that Gauss gives without proof towards the end of the fifth section, it appears that Gauss already knew the \href{https://en.wikipedia.org/wiki/Class_number_formula}{class number formula} in 1801.

In the last section, Gauss gives proof for the \href{https://en.wikipedia.org/wiki/Constructible_polygon}{constructibility} of a regular \href{https://en.wikipedia.org/wiki/Heptadecagon}{heptadecagon} (17-sided polygon) with \href{https://en.wikipedia.org/wiki/Compass_and_straightedge_constructions}{straightedge and compass} by reducing this geometrical problem to an algebraic one. This was the first progress in regular polygon construction in over 2000 years. He shows that a regular polygon is constructible if the number of its sides is either a \href{https://en.wikipedia.org/wiki/Power_of_2}{power of 2} or the product of a power of 2 and any number of distinct \href{https://en.wikipedia.org/wiki/Fermat_prime}{Fermat primes}. In the same section, he gives a result on the number of solutions of certain cubic polynomials with coefficients in \href{https://en.wikipedia.org/wiki/Finite_field}{finite fields}, which amounts to counting integral points on an \href{https://en.wikipedia.org/wiki/Elliptic_curve}{elliptic curve}. An unfinished chapter, consisting of work done during 1797--1799, was found among his papers after his death.


\paragraph{Further investigations}\label{Further_investigations}

One of Gauss\textquotesingle s first results was the empirically found conjecture of 1792 -- the later called \href{https://en.wikipedia.org/wiki/Prime_number_theorem}{prime number theorem} -- giving an estimation of the number of prime numbers by using the \href{https://en.wikipedia.org/wiki/Logarithmic_integral_function}{integral logarithm}.

In 1816, \href{https://en.wikipedia.org/wiki/Heinrich_Wilhelm_Matthias_Olbers}{Olbers} encouraged Gauss to compete for a prize from the French Academy for a proof for \href{https://en.wikipedia.org/wiki/Fermat's_Last_Theorem}{Fermat\textquotesingle s Last Theorem}; he refused, considering the topic uninteresting. However, after his death a short undated paper was found with proofs of the theorem for the cases \emph{n} = 3 and \emph{n} = 5. The particular case of \emph{n} = 3 was proved much earlier by \href{https://en.wikipedia.org/wiki/Leonhard_Euler}{Leonhard Euler}, but Gauss developed a more streamlined proof which made use of \href{https://en.wikipedia.org/wiki/Eisenstein_integers}{Eisenstein integers}; though more general, the proof was simpler than in the real integers case.

Gauss contributed to solving the \href{https://en.wikipedia.org/wiki/Kepler_conjecture}{Kepler conjecture} in 1831 with the proof that a \href{https://en.wikipedia.org/wiki/Close-packing_of_equal_spheres}{greatest packing density} of spheres in the three-dimensional space is given when the centres of the spheres form a \href{https://en.wikipedia.org/wiki/Cubic_crystal_system}{cubic face-centred} arrangement, when he reviewed a book of \href{https://en.wikipedia.org/wiki/Ludwig_August_Seeber}{Ludwig August Seeber} on the theory of reduction of positive ternary quadratic forms. Having noticed some lacks in Seeber\textquotesingle s proof, he simplified many of his arguments, proved the central conjecture, and remarked that this theorem is equivalent to the Kepler conjecture for regular arrangements.

In two papers on \href{https://en.wikipedia.org/wiki/Quartic_reciprocity}{biquadratic residues} (1828, 1832) Gauss introduced the \href{https://en.wikipedia.org/wiki/Ring_theory}{ring} of \href{https://en.wikipedia.org/wiki/Gaussian_integers}{Gaussian integers} {{\(\mathbb{Z}\lbrack i\rbrack\)}{}}, showed that it is a \href{https://en.wikipedia.org/wiki/Unique_factorization_domain}{unique factorization domain}, and generalized some key arithmetic concepts, such as \href{https://en.wikipedia.org/wiki/Fermat's_little_theorem}{Fermat\textquotesingle s little theorem} and \href{https://en.wikipedia.org/wiki/Gauss's_lemma_(number_theory)}{Gauss\textquotesingle s lemma}. The main objective of introducing this ring was to formulate the law of biquadratic reciprocity -- as Gauss discovered, rings of complex integers are the natural setting for such higher reciprocity laws.

In the second paper, he stated the general law of biquadratic reciprocity and proved several special cases of it. In an earlier publication from 1818 containing his fifth and sixth proofs of quadratic reciprocity, he claimed the techniques of these proofs (\href{https://en.wikipedia.org/wiki/Gauss_sum}{Gauss sums}) can be applied to prove higher reciprocity laws.


\subsubsection{Analysis}\label{Analysis}

One of Gauss\textquotesingle s first discoveries was the notion of the \href{https://en.wikipedia.org/wiki/Arithmetic-geometric_mean}{arithmetic-geometric mean} (AGM) of two positive real numbers. He discovered its relation to elliptic integrals in the years 1798--1799 through \href{https://en.wikipedia.org/wiki/Landen's_transformation}{Landen\textquotesingle s transformation}, and a diary entry recorded the discovery of the connection of \href{https://en.wikipedia.org/wiki/Gauss's_constant}{Gauss\textquotesingle s constant} to \href{https://en.wikipedia.org/wiki/Lemniscatic_elliptic_functions}{lemniscatic elliptic functions}, a result that Gauss stated "will surely open an entirely new field of analysis". He also made early inroads into the more formal issues of the foundations of \href{https://en.wikipedia.org/wiki/Complex_analysis}{complex analysis}, and from a letter to Bessel in 1811 it is clear that he knew the "fundamental theorem of complex analysis" -- \href{https://en.wikipedia.org/wiki/Cauchy's_integral_theorem}{Cauchy\textquotesingle s integral theorem} -- and understood the notion of \href{https://en.wikipedia.org/wiki/Residue_(complex_analysis)}{complex residues} when integrating around \href{https://en.wikipedia.org/wiki/Pole_(complex_analysis)}{poles}.

\href{https://en.wikipedia.org/wiki/Pentagonal_number_theorem}{Euler\textquotesingle s pentagonal numbers theorem}, together with other researches on the AGM and lemniscatic functions, led him to plenty of results on \href{https://en.wikipedia.org/wiki/Jacobi_theta_functions}{Jacobi theta functions}, culminating in the discovery in 1808 of the later called \href{https://en.wikipedia.org/wiki/Jacobi_triple_product_identity}{Jacobi triple product identity}, which includes Euler\textquotesingle s theorem as a special case. His works show that he knew modular transformations of order 3, 5, 7 for elliptic functions since 1808.

Several mathematical fragments in his \href{https://en.wikipedia.org/wiki/Nachlass}{Nachlass} indicate that he knew parts of the modern theory of \href{https://en.wikipedia.org/wiki/Modular_forms}{modular forms}. In his work on the \href{https://en.wikipedia.org/wiki/Multivalued_function}{multivalued} AGM of two complex numbers, he discovered a deep connection between the infinitely many values of the AGM and its two "simplest values". In his unpublished writings he recognized and made a sketch of the key concept of \href{https://en.wikipedia.org/wiki/Fundamental_domain}{fundamental domain} for the \href{https://en.wikipedia.org/wiki/Modular_group}{modular group}. One of Gauss\textquotesingle s sketches of this kind was a drawing of a \href{https://en.wikipedia.org/wiki/Tessellation}{tessellation} of the \href{https://en.wikipedia.org/wiki/Unit_disk}{unit disk} by "equilateral" \href{https://en.wikipedia.org/wiki/Hyperbolic_triangle}{hyperbolic triangles} with all angles equal to {{\(\pi/4\)}{}}.

An example of Gauss\textquotesingle s insight in analysis is the cryptic remark that the principles of circle division by compass and straightedge can also be applied to the division of the \href{https://en.wikipedia.org/wiki/Lemniscate_curve}{lemniscate curve}, which inspired Abel\textquotesingle s theorem on lemniscate division. Another example is his publication "Summatio quarundam serierum singularium" (1811) on the determination of the sign of \href{https://en.wikipedia.org/wiki/Quadratic_Gauss_sums?action=edit&redlink=1}{quadratic Gauss sums}, in which he solved the main problem by introducing \href{https://en.wikipedia.org/wiki/Gaussian_binomial_coefficient}{q-analogs of binomial coefficients} and manipulating them by several original identities that seem to stem from his work on elliptic function theory; however, Gauss cast his argument in a formal way that does not reveal its origin in elliptic function theory, and only the later work of mathematicians such as \href{https://en.wikipedia.org/wiki/Carl_Gustav_Jacob_Jacobi}{Jacobi} and \href{https://en.wikipedia.org/wiki/Hermite}{Hermite} has exposed the crux of his argument.

In the "Disquisitiones generales circa series infinitam..." (1813), he provides the first systematic treatment of the general \href{https://en.wikipedia.org/wiki/Hypergeometric_function}{hypergeometric function} {{\(F(\alpha,\beta,\gamma,x)\)}{}}, and shows that many of the functions known at the time are special cases of the hypergeometric function. This work is the first exact inquiry into \href{https://en.wikipedia.org/wiki/Convergent_series}{convergence} of infinite series in the history of mathematics. Furthermore, it deals with infinite \href{https://en.wikipedia.org/wiki/Continued_fraction}{continued fractions} arising as ratios of hypergeometric functions, which are now called \href{https://en.wikipedia.org/wiki/Gauss_continued_fraction}{Gauss continued fractions}.

In 1823, Gauss won the prize of the Danish Society with an essay on \href{https://en.wikipedia.org/wiki/Conformal_mapping}{conformal mappings}, which contains several developments that pertain to the field of complex analysis. Gauss stated that angle-preserving mappings in the complex plane must be complex \href{https://en.wikipedia.org/wiki/Analytic_function}{analytic functions}, and used the later-named \href{https://en.wikipedia.org/wiki/Beltrami_equation}{Beltrami equation} to prove the existence of \href{https://en.wikipedia.org/wiki/Isothermal_coordinates}{isothermal coordinates} on analytic surfaces. The essay concludes with examples of conformal mappings into a sphere and an \href{https://en.wikipedia.org/wiki/Ellipsoid_of_revolution}{ellipsoid of revolution}.


\paragraph{Numerical analysis}\label{Numerical_analysis}

Gauss often deduced theorems \href{https://en.wikipedia.org/wiki/Inductive_reasoning}{inductively} from numerical data he had collected empirically. As such, the use of efficient algorithms to facilitate calculations was vital to his research, and he made many contributions to \href{https://en.wikipedia.org/wiki/Numerical_analysis}{numerical analysis}, such as the method of \href{https://en.wikipedia.org/wiki/Gaussian_quadrature}{Gaussian quadrature}, published in 1816.

In a private letter to \href{https://en.wikipedia.org/wiki/Christian_Ludwig_Gerling}{Gerling} from 1823, he described a solution of a 4x4 system of linear equations with the \href{https://en.wikipedia.org/wiki/Gauss-Seidel_method}{Gauss-Seidel method} -- an "indirect" \href{https://en.wikipedia.org/wiki/Iterative_method}{iterative method} for the solution of linear systems, and recommended it over the usual method of "direct elimination" for systems of more than two equations.

Gauss invented an algorithm for calculating what is now called \href{https://en.wikipedia.org/wiki/Discrete_Fourier_transform}{discrete Fourier transforms} when calculating the orbits of Pallas and Juno in 1805, 160 years before \href{https://en.wikipedia.org/wiki/James_Cooley}{Cooley} and \href{https://en.wikipedia.org/wiki/John_Tukey}{Tukey} found their similar \href{https://en.wikipedia.org/wiki/Cooley-Tukey_FFT_algorithm}{Cooley--Tukey algorithm}. He developed it as a \href{https://en.wikipedia.org/wiki/Trigonometric_interpolation}{trigonometric interpolation} method, but the paper \emph{Theoria Interpolationis Methodo Nova Tractata} was published only posthumously in 1876, well after \href{https://en.wikipedia.org/wiki/Joseph_Fourier}{Joseph Fourier}\textquotesingle s introduction of the subject in 1807.


\subsubsection{Geometry}\label{Geometry}


\paragraph{Differential geometry}\label{Differential_geometry}

The geodetic survey of \href{https://en.wikipedia.org/wiki/Hanover}{Hanover} fuelled Gauss\textquotesingle s interest in \href{https://en.wikipedia.org/wiki/Differential_geometry}{differential geometry} and \href{https://en.wikipedia.org/wiki/Topology}{topology}, fields of mathematics dealing with \href{https://en.wikipedia.org/wiki/Curve}{curves} and \href{https://en.wikipedia.org/wiki/Surface_(topology)}{surfaces}. This led him in 1828 to the publication of a work that marks the birth of modern \href{https://en.wikipedia.org/wiki/Differential_geometry_of_surfaces}{differential geometry of surfaces}, as it departed from the traditional ways of treating surfaces as \href{https://en.wikipedia.org/wiki/Cartesian_coordinate_system}{cartesian graphs} of functions of two variables, and that initiated the exploration of surfaces from the "inner" point of view of a two-dimensional being constrained to move on it. As a result, the \href{https://en.wikipedia.org/wiki/Theorema_Egregium}{Theorema Egregium} (\emph{remarkable theorem}), established a property of the notion of \href{https://en.wikipedia.org/wiki/Gaussian_curvature}{Gaussian curvature}. Informally, the theorem says that the curvature of a surface can be determined entirely by measuring \href{https://en.wikipedia.org/wiki/Angle}{angles} and \href{https://en.wikipedia.org/wiki/Distance}{distances} on the surface, regardless of the \href{https://en.wikipedia.org/wiki/Embedding}{embedding} of the surface in three-dimensional or two-dimensional space.

The Theorema Egregium leads to the abstraction of surfaces as doubly-extended \href{https://en.wikipedia.org/wiki/Manifold}{manifolds}; it clarifies the distinction between the intrinsic properties of the manifold (the \href{https://en.wikipedia.org/wiki/Metric_tensor}{metric}) and its physical realization in ambient space. A consequence is the impossibility of an isometric transformation between surfaces of different Gaussian curvature. This means practically that a \href{https://en.wikipedia.org/wiki/Sphere}{sphere} or an \href{https://en.wikipedia.org/wiki/Ellipsoid}{ellipsoid} cannot be transformed to a plane without distortion, which causes a fundamental problem in designing \href{https://en.wikipedia.org/wiki/Map_projection}{projections} for geographical maps. A portion of this essay is dedicated to a profound study of \href{https://en.wikipedia.org/wiki/Geodesic}{geodesics}. In particular, Gauss proves the local \href{https://en.wikipedia.org/wiki/Gauss–Bonnet_theorem}{Gauss--Bonnet theorem} on geodesic triangles, and generalizes \href{https://en.wikipedia.org/wiki/Legendre's_theorem_on_spherical_triangles}{Legendre\textquotesingle s theorem on spherical triangles} to geodesic triangles on arbitrary surfaces with continuous curvature; he found that the angles of a "sufficiently small" geodesic triangle deviate from that of a planar triangle of the same sides in a way that depends only on the values of the surface curvature at the vertices of the triangle, regardless of the behaviour of the surface in the triangle interior.

Gauss\textquotesingle s memoir from 1828 lacks the conception of \href{https://en.wikipedia.org/wiki/Geodesic_curvature}{geodesic curvature}. However, in a previously unpublished manuscript, very likely written in 1822--1825, he introduced the term "side curvature" (German: "Seitenkrümmung") and proved its \href{https://en.wikipedia.org/wiki/Invariant_(mathematics)}{invariance} under isometric transformations, a result that was later obtained by \href{https://en.wikipedia.org/wiki/Ferdinand_Minding}{Ferdinand Minding} and published by him in 1830. This Gauss paper contains the core of his lemma on total curvature, but also its generalization, found and proved by \href{https://en.wikipedia.org/wiki/Pierre_Ossian_Bonnet}{Pierre Ossian Bonnet} in 1848 and known as the \href{https://en.wikipedia.org/wiki/Gauss–Bonnet_theorem}{Gauss--Bonnet theorem}.


\paragraph{Non-Euclidean geometry}\label{Non-Euclidean_geometry}

\begin{figure}
\centering
\href{https://en.wikipedia.org/wiki/File:Bendixen_-_Carl_Friedrich_Gauß,_1828.jpg}{\includegraphics[width=1.97917in,height=2.25in]{250px-Bendixen_-_Carl_Friedrich_Gau_1828.jpg}}
\caption{Lithograph of Gauss by \href{https://en.wikipedia.org/wiki/Siegfried_Bendixen}{Siegfried Bendixen} (1828)}\label{mwBZY}
\end{figure}

During Gauss\textquotesingle s lifetime, the \href{https://en.wikipedia.org/wiki/Parallel_postulate}{parallel postulate} of \href{https://en.wikipedia.org/wiki/Euclidean_geometry}{Euclidean geometry} was heavily discussed. Numerous efforts were made to prove it in the frame of the Euclidean \href{https://en.wikipedia.org/wiki/Axiom}{axioms}, whereas some mathematicians discussed the possibility of geometrical systems without it. Gauss thought about the basics of geometry from the 1790s on, but only realized in the 1810s that a non-Euclidean geometry without the parallel postulate could solve the problem. In a letter to \href{https://en.wikipedia.org/wiki/Franz_Taurinus}{Franz Taurinus} of 1824, he presented a short comprehensible outline of what he named a "\href{https://en.wikipedia.org/wiki/Non-Euclidean_geometry}{non-Euclidean geometry}", but he strongly forbade Taurinus to make any use of it. Gauss is credited with having been the one to first discover and study non-Euclidean geometry, even coining the term as well.

The first publications on non-Euclidean geometry in the history of mathematics were authored by \href{https://en.wikipedia.org/wiki/Nikolai_Lobachevsky}{Nikolai Lobachevsky} in 1829 and \href{https://en.wikipedia.org/wiki/Janos_Bolyai}{Janos Bolyai} in 1832. In the following years, Gauss wrote his ideas on the topic but did not publish them, thus avoiding influencing the contemporary scientific discussion. Gauss commended the ideas of Janos Bolyai in a letter to his father and university friend Farkas Bolyai claiming that these were congruent to his own thoughts of some decades. However, it is not quite clear to what extent he preceded Lobachevsky and Bolyai, as his written remarks are vague and obscure.

\href{https://en.wikipedia.org/wiki/Wolfgang_Sartorius_von_Waltershausen}{Sartorius} first mentioned Gauss\textquotesingle s work on non-Euclidean geometry in 1856, but only the publication of Gauss\textquotesingle s \href{https://en.wikipedia.org/wiki/Nachlass}{Nachlass} in Volume VIII of the Collected Works (1900) showed Gauss\textquotesingle s ideas on the matter, at a time when non-Euclidean geometry was still an object of some controversy.


\paragraph{Early topology}\label{Early_topology}

Gauss was also an early pioneer of \href{https://en.wikipedia.org/wiki/Topology}{topology} or \emph{Geometria Situs}, as it was called in his lifetime. The first proof of the \href{https://en.wikipedia.org/wiki/Fundamental_theorem_of_algebra}{fundamental theorem of algebra} in 1799 contained an essentially topological argument; fifty years later, he further developed the topological argument in his fourth proof of this theorem.

\begin{figure}
\centering
\href{https://en.wikipedia.org/wiki/File:Carl_Friedrich_Gauß,_Büste_von_Heinrich_Hesemann,_1855.jpg}{\includegraphics[width=1.97917in,height=2.875in]{250px-Carl_Friedrich_Gau_B_ste_von_Heinrich_Hesemann_1855.jpg}}
\caption{Gauss bust by Heinrich Hesemann (1855)}\label{mwBfQ}
\end{figure}

Another encounter with topological notions occurred to him in the course of his astronomical work in 1804, when he determined the limits of the region on the \href{https://en.wikipedia.org/wiki/Celestial_sphere}{celestial sphere} in which comets and asteroids might appear, and which he termed "Zodiacus". He discovered that if the Earth\textquotesingle s and comet\textquotesingle s orbits are \href{https://en.wikipedia.org/wiki/Link_(knot_theory)}{linked}, then by topological reasons the Zodiacus is the entire sphere. In 1848, in the context of the discovery of the asteroid \href{https://en.wikipedia.org/wiki/7_Iris}{7 Iris}, he published a further qualitative discussion of the Zodiacus.

In Gauss\textquotesingle s letters of 1820--1830, he thought intensively on topics with close affinity to Geometria Situs, and became gradually conscious of semantic difficulty in this field. Fragments from this period reveal that he tried to classify "tract figures", which are closed plane curves with a finite number of transverse self-intersections, that may also be planar projections of \href{https://en.wikipedia.org/wiki/Knot_theory}{knots}. To do so he devised a symbolical scheme, the \href{https://en.wikipedia.org/wiki/Gauss_code}{Gauss code}, that in a sense captured the characteristic features of tract figures.

In a fragment from 1833, Gauss defined the \href{https://en.wikipedia.org/wiki/Linking_number}{linking number} of two space curves by a certain double integral, and in doing so provided for the first time an analytical formulation of a topological phenomenon. On the same note, he lamented the little progress made in Geometria Situs, and remarked that one of its central problems will be "to count the intertwinings of two closed or infinite curves". His notebooks from that period reveal that he was also thinking about other topological objects such as \href{https://en.wikipedia.org/wiki/Braid_theory}{braids} and \href{https://en.wikipedia.org/wiki/Tangle_(mathematics)}{tangles}.

Gauss\textquotesingle s influence in later years to the emerging field of topology, which he held in high esteem, was through occasional remarks and oral communications to Mobius and Listing.


\subsubsection{Minor mathematical accomplishments}\label{Minor_mathematical_accomplishments}

Gauss applied the concept of complex numbers to solve well-known problems in a new concise way. For example, in a short note from 1836 on geometric aspects of the ternary forms and their application to crystallography, he stated the \href{https://en.wikipedia.org/wiki/Axonometry}{fundamental theorem of axonometry}, which tells how to represent a 3D cube on a 2D plane with complete accuracy, via complex numbers. He described rotations of this sphere as the action of certain \href{https://en.wikipedia.org/wiki/Mobius_transformation}{linear fractional transformations} on the extended complex plane, and gave a proof for the geometric theorem that the \href{https://en.wikipedia.org/wiki/Altitude_(triangle)}{altitudes} of a triangle always meet in a single \href{https://en.wikipedia.org/wiki/Orthocenter}{orthocenter}.

Gauss was concerned with \href{https://en.wikipedia.org/wiki/John_Napier}{John Napier}\textquotesingle s "\href{https://en.wikipedia.org/wiki/Pentagramma_mirificum}{Pentagramma mirificum}" -- a certain spherical \href{https://en.wikipedia.org/wiki/Pentagram}{pentagram} -- for several decades; he approached it from various points of view, and gradually gained a full understanding of its geometric, algebraic, and analytic aspects. In particular, in 1843 he stated and proved several theorems connecting elliptic functions, Napier spherical pentagons, and Poncelet pentagons in the plane.

Furthermore, he contributed a solution to the problem of constructing the largest-area ellipse inside a given \href{https://en.wikipedia.org/wiki/Quadrilateral}{quadrilateral}, and discovered a surprising result about the computation of area of \href{https://en.wikipedia.org/wiki/Pentagon}{pentagons}.


\subsection{Sciences}\label{Sciences}


\subsubsection{Astronomy}\label{Astronomy}

\begin{figure}
\centering
\caption{Carl Friedrich Gauss 1803 by Johann Christian August Schwartz}\label{mwBl0}
\end{figure}

On 1 January 1801, Italian astronomer \href{https://en.wikipedia.org/wiki/Giuseppe_Piazzi}{Giuseppe Piazzi} discovered a new celestial object, presumed it to be the long searched planet between Mars and Jupiter according to the so-called \href{https://en.wikipedia.org/wiki/Titius–Bode_law}{Titius--Bode law}, and named it \href{https://en.wikipedia.org/wiki/Ceres_(dwarf_planet)}{Ceres}. He could track it only for a short time until it disappeared behind the glare of the Sun. The mathematical tools of the time were not sufficient to predict the location of its reappearance from the few data available. Gauss tackled the problem and predicted a position for possible rediscovery in December 1801. This turned out to be accurate within a half-degree when \href{https://en.wikipedia.org/wiki/Franz_Xaver_von_Zach}{Franz Xaver von Zach} on 7 and 31 December at \href{https://en.wikipedia.org/wiki/Gotha_Observatory}{Gotha}, and independently \href{https://en.wikipedia.org/wiki/Heinrich_Wilhelm_Matthäus_Olbers}{Heinrich Olbers} on 1 and 2 January in \href{https://en.wikipedia.org/wiki/Bremen}{Bremen}, identified the object near the predicted position.

\href{https://en.wikipedia.org/wiki/Gauss's_method}{Gauss\textquotesingle s method} leads to an equation of the eighth degree, of which one solution, the Earth\textquotesingle s orbit, is known. The solution sought is then separated from the remaining six based on physical conditions. In this work, Gauss used comprehensive approximation methods which he created for that purpose.

The discovery of Ceres led Gauss to the theory of the motion of planetoids disturbed by large planets, eventually published in 1809 as \emph{Theoria motus corporum coelestium in sectionibus conicis solem ambientum}. It introduced the \href{https://en.wikipedia.org/wiki/Gaussian_gravitational_constant}{Gaussian gravitational constant}.

Since the new asteroids had been discovered, Gauss occupied himself with the \href{https://en.wikipedia.org/wiki/Perturbation_(astronomy)}{perturbations} of their \href{https://en.wikipedia.org/wiki/Orbital_elements}{orbital elements}. Firstly he examined Ceres with analytical methods similar to those of Laplace, but his favorite object was \href{https://en.wikipedia.org/wiki/(2)_Pallas}{Pallas}, because of its great \href{https://en.wikipedia.org/wiki/Eccentricity_(astronomy)}{eccentricity} and \href{https://en.wikipedia.org/wiki/Orbital_inclination}{orbital inclination}, whereby Laplace\textquotesingle s method did not work. Gauss used his own tools: the \href{https://en.wikipedia.org/wiki/Arithmetic–geometric_mean}{arithmetic--geometric mean}, the \href{https://en.wikipedia.org/wiki/Hypergeometric_function}{hypergeometric function}, and his method of interpolation. He found an \href{https://en.wikipedia.org/wiki/Orbital_resonance\#Coincidental_'near'_ratios_of_mean_motion}{orbital resonance} with \href{https://en.wikipedia.org/wiki/Jupiter}{Jupiter} in proportion 18:7 in 1812; Gauss gave this result as \href{https://en.wikipedia.org/wiki/Cipher}{cipher}, and gave the explicit meaning only in letters to Olbers and Bessel. After long years of work, he finished it in 1816 without a result that seemed sufficient to him. This marked the end of his activities in theoretical astronomy.

\begin{figure}
\centering
\href{https://en.wikipedia.org/wiki/File:Goettingen_Sternwarte_Besemann.png}{\includegraphics[width=2.60417in,height=1.65625in]{250px-Goettingen_Sternwarte_Besemann.png}}
\caption{Göttingen observatory seen from the North-west (by Friedrich Besemann, {c.}{{ }1835})}\label{mwBqU}
\end{figure}

One fruit of Gauss\textquotesingle s research on Pallas perturbations was the \emph{Determinatio Attractionis...} (1818) on a method of theoretical astronomy that later became known as the "elliptic ring method". It introduced an averaging conception in which a planet in orbit is replaced by a fictitious ring with mass density proportional to the time the planet takes to follow the corresponding orbital arcs. Gauss presents the method of evaluating the gravitational attraction of such an elliptic ring, which includes several steps; one of them involves a direct application of the arithmetic-geometric mean (AGM) algorithm to calculate an \href{https://en.wikipedia.org/wiki/Elliptic_integral}{elliptic integral}.

Even after Gauss\textquotesingle s contributions to theoretical astronomy came to an end, more practical activities in \href{https://en.wikipedia.org/wiki/Observational_astronomy}{observational astronomy} continued and occupied him during his entire career. As early as 1799, Gauss dealt with the determination of longitude by use of the lunar parallax, for which he developed more convenient formulas than those were in common use. After appointment as director of observatory he attached importance to the fundamental astronomical constants in correspondence with Bessel. Gauss himself provided tables of \href{https://en.wikipedia.org/wiki/Astronomical_nutation}{nutation} and \href{https://en.wikipedia.org/wiki/Aberration_(astronomy)}{aberration}, solar coordinates, and refraction. He made many contributions to \href{https://en.wikipedia.org/wiki/Spherical_geometry}{spherical geometry}, and in this context solved some practical problems about \href{https://en.wikipedia.org/wiki/Celestial_navigation}{navigation by stars}. He published a great number of observations, mainly on minor planets and comets; his last observation was the \href{https://en.wikipedia.org/wiki/Solar_eclipse_of_28_July_1851?action=edit&redlink=1}{solar eclipse of 28 July 1851}.


\subsubsection{Chronology}\label{Chronology}

Gauss\textquotesingle s first publication following his doctoral thesis dealt with the determination of the \href{https://en.wikipedia.org/wiki/Date_of_Easter\#Gauss's_Easter_algorithm}{date of Easter} (1800), an elementary mathematical topic. Gauss aimed to present a convenient algorithm for people without any knowledge of ecclesiastical or even astronomical chronology, and thus avoided the usual terms of \href{https://en.wikipedia.org/wiki/Golden_number_(time)}{golden number}, \href{https://en.wikipedia.org/wiki/Epact}{epact}, \href{https://en.wikipedia.org/wiki/Solar_cycle_(calendar)}{solar cycle}, \href{https://en.wikipedia.org/wiki/Dominical_letter}{domenical letter}, and any religious connotations. This choice of topic likely had historical grounds. The replacement of the \href{https://en.wikipedia.org/wiki/Julian_calendar}{Julian calendar} by the \href{https://en.wikipedia.org/wiki/Gregorian_calendar}{Gregorian calendar} had caused confusion in the \href{https://en.wikipedia.org/wiki/Holy_Roman_Empire}{Holy Roman Empire} since the 16th century and was not finished in Germany until 1700, when the difference of eleven days was deleted. Even after this, Easter fell on different dates in Protestant and Catholic territories, until this difference was abolished by agreement in 1776. In the Protestant states, such as the Duchy of Brunswick, the Easter of 1777, five weeks before Gauss\textquotesingle s birth, was the first one calculated in the new manner.


\subsubsection{Error theory}\label{Error_theory}

Gauss likely used the \href{https://en.wikipedia.org/wiki/Least_squares}{method of least squares} to minimize the impact of \href{https://en.wikipedia.org/wiki/Observational_error}{measurement error} when calculating the orbit of Ceres. The method was published first by \href{https://en.wikipedia.org/wiki/Adrien-Marie_Legendre}{Adrien-Marie Legendre} in 1805, but Gauss claimed in \emph{Theoria motus} (1809) that he had been using it since 1794 or 1795. In the history of statistics, this disagreement is called the "priority dispute over the discovery of the method of least squares". Gauss proved that the method has the lowest sampling variance within the class of linear unbiased estimators under the assumption of \href{https://en.wikipedia.org/wiki/Normal_distribution}{normally distributed} errors (\href{https://en.wikipedia.org/wiki/Gauss–Markov_theorem}{Gauss--Markov theorem}), in the two-part paper \emph{Theoria combinationis observationum erroribus minimis obnoxiae} (1823).

In the first paper he proved \href{https://en.wikipedia.org/wiki/Gauss's_inequality}{Gauss\textquotesingle s inequality} (a \href{https://en.wikipedia.org/wiki/Chebyshev's_inequality}{Chebyshev-type inequality}) for \href{https://en.wikipedia.org/wiki/Unimodality}{unimodal distributions}, and stated without proof another inequality for \href{https://en.wikipedia.org/wiki/Moment_(mathematics)}{moments} of the fourth order (a special case of the Gauss-Winckler inequality). He derived lower and upper bounds for the \href{https://en.wikipedia.org/wiki/Variance}{variance} of the \href{https://en.wikipedia.org/wiki/Sample_variance}{sample variance}. In the second paper, Gauss described \href{https://en.wikipedia.org/wiki/Recursive_least_squares_filter}{recursive least squares methods}, which he discovered. Gauss\textquotesingle s work on the theory of errors was extended in several directions by the geodesist \href{https://en.wikipedia.org/wiki/Friedrich_Robert_Helmert}{Friedrich Robert Helmert} to the \href{https://en.wikipedia.org/wiki/Gauss-Helmert_model}{Gauss-Helmert model}.

Gauss also contributed to problems in \href{https://en.wikipedia.org/wiki/Probability_theory}{probability theory} that are not directly concerned with the theory of errors. One example appears as a diary note where he tried to describe the asymptotic distribution of entries in the continued fraction expansion of a random number uniformly distributed in \emph{(0,1)}. He derived this distribution, now known as the \href{https://en.wikipedia.org/wiki/Gauss-Kuzmin_distribution}{Gauss-Kuzmin distribution}, as a by-product of the discovery of the \href{https://en.wikipedia.org/wiki/Ergodicity}{ergodicity} of the \href{https://en.wikipedia.org/wiki/Gauss-Kuzmin-Wirsing_operator}{Gauss map for continued fractions}. Gauss\textquotesingle s solution is the first-ever result in the metrical theory of continued fractions.


\subsubsection{Geodesy}\label{Geodesy}

\begin{figure}
\centering
\href{https://en.wikipedia.org/wiki/File:Georg_IV_Erlass_Landvermessung.jpg}{\includegraphics[width=1.97917in,height=3.04167in]{250px-Georg_IV_Erlass_Landvermessung.jpg}}
\caption{Order of King \href{https://en.wikipedia.org/wiki/George_IV}{George IV} from 9 May 1820 to the triangulation project (with the additional signature of Count \href{https://en.wikipedia.org/wiki/Ernst_zu_Münster}{Ernst zu Münster} below)}\label{mwByE}
\end{figure}

\begin{figure}
\centering
\href{https://en.wikipedia.org/wiki/File:Gauß_Heliotrop@20170430.JPG}{\includegraphics[width=2.60417in,height=1.73958in]{250px-Gau_Heliotrop_20170430.JPG}}
\caption{The \href{https://en.wikipedia.org/wiki/Heliotrope_(instrument)}{heliotrope}}\label{mwByc}
\end{figure}

\begin{figure}
\centering
\href{https://en.wikipedia.org/wiki/File:Vize-Heliotrop_Gauß-Ausstellung_Bomann-Museum_(1).jpg}{\includegraphics[width=2.60417in,height=2.125in]{250px-Vize-Heliotrop_Gau_-Ausstellung_Bomann-Museum_1_.jpg}}
\caption{Gauss\textquotesingle s vice heliotrope, a \href{https://en.wikipedia.org/wiki/Edward_Troughton}{Troughton} sextant with additional mirror}\label{mwByw}
\end{figure}

Gauss was busy with geodetic problems since 1799 when he helped \href{https://en.wikipedia.org/wiki/Karl_Ludwig_von_Lecoq}{Karl Ludwig von Lecoq} with calculations during his \href{https://en.wikipedia.org/wiki/Surveying}{survey} in \href{https://en.wikipedia.org/wiki/Westphalia}{Westphalia}. Beginning in 1804, he taught himself some practical geodesy in Brunswick and Göttingen.

Since 1816, Gauss\textquotesingle s former student \href{https://en.wikipedia.org/wiki/Heinrich_Christian_Schumacher}{Heinrich Christian Schumacher}, then professor in \href{https://en.wikipedia.org/wiki/Copenhagen}{Copenhagen}, but living in \href{https://en.wikipedia.org/wiki/Altona,_Hamburg}{Altona} (\href{https://en.wikipedia.org/wiki/Duchy_of_Holstein}{Holstein}) near \href{https://en.wikipedia.org/wiki/Hamburg}{Hamburg} as head of an observatory, carried out a \href{https://en.wikipedia.org/wiki/Triangulation}{triangulation} of the \href{https://en.wikipedia.org/wiki/Jutland}{Jutland} peninsula from \href{https://en.wikipedia.org/wiki/Skagen}{Skagen} in the north to \href{https://en.wikipedia.org/wiki/Lauenburg}{Lauenburg} in the south. This project was the basis for map production but also aimed at determining the geodetic arc between the terminal sites. Data from geodetic arcs were used to determine the dimensions of the earth \href{https://en.wikipedia.org/wiki/Geoid}{geoid}, and long arc distances brought more precise results. Schumacher asked Gauss to continue this work further to the south in the Kingdom of Hanover; Gauss agreed after a short time of hesitation. Finally, in May 1820, King \href{https://en.wikipedia.org/wiki/George_IV}{George IV} gave the order to Gauss.

An \href{https://en.wikipedia.org/wiki/Arc_measurement}{arc measurement} needs a precise astronomical determination of at least two points in the \href{https://en.wikipedia.org/wiki/Triangulation_network}{network}. Gauss and Schumacher used the coincidence that both observatories in Göttingen and Altona, in the garden of Schumacher\textquotesingle s house, lay nearly in the same \href{https://en.wikipedia.org/wiki/Longitude}{longitude}. The \href{https://en.wikipedia.org/wiki/Latitude}{latitude} was measured with both their instruments and a \href{https://en.wikipedia.org/wiki/Zenith_sector}{zenith sector} of \href{https://en.wikipedia.org/wiki/Jesse_Ramsden}{Ramsden} that was transported to both observatories.

Gauss and Schumacher had already determined some angles between \href{https://en.wikipedia.org/wiki/Lüneburg}{Lüneburg}, Hamburg, and Lauenburg for the geodetic connection in October 1818. During the summers of 1821 until 1825 Gauss directed the triangulation work personally, from \href{https://en.wikipedia.org/wiki/Thuringia}{Thuringia} in the south to the river \href{https://en.wikipedia.org/wiki/Elbe}{Elbe} in the north. The \href{https://en.wikipedia.org/wiki/Triangle}{triangle} between \href{https://en.wikipedia.org/wiki/Hoher_Hagen_(Dransfeld)}{Hoher Hagen}, \href{https://en.wikipedia.org/wiki/Großer_Inselsberg}{Großer Inselsberg} in the \href{https://en.wikipedia.org/wiki/Thuringian_Forest}{Thuringian Forest}, and \href{https://en.wikipedia.org/wiki/Brocken}{Brocken} in the \href{https://en.wikipedia.org/wiki/Harz}{Harz} mountains was the largest one Gauss had ever measured with a maximum size of {107}{~}{km (66.5 miles)}. In the thinly populated \href{https://en.wikipedia.org/wiki/Lüneburg_Heath}{Lüneburg Heath} without significant natural summits or artificial buildings, he had difficulties finding suitable triangulation points; sometimes cutting lanes through the vegetation was necessary.

For pointing signals, Gauss invented a new instrument with movable mirrors and a small telescope that reflects the sunbeams to the triangulation points, and named it \emph{\href{https://en.wikipedia.org/wiki/Heliotrope_(instrument)}{heliotrope}}. Another suitable construction for the same purpose was a \href{https://en.wikipedia.org/wiki/Sextant}{sextant} with an additional mirror which he named \emph{vice heliotrope}. Gauss was assisted by soldiers of the Hanoverian army, among them his eldest son Joseph. Gauss took part in the \href{https://en.wikipedia.org/wiki/Baseline_(surveying)}{baseline} measurement (\href{https://en.wikipedia.org/wiki/Braak_Base_Line}{Braak Base Line}) of Schumacher in the village of \href{https://en.wikipedia.org/wiki/Braak,_Schleswig-Holstein}{Braak} near Hamburg in 1820, and used the result for the evaluation of the Hanoverian triangulation.

An additional result was a better value for the \href{https://en.wikipedia.org/wiki/Flattening}{flattening} of the approximative \href{https://en.wikipedia.org/wiki/Earth_ellipsoid}{Earth ellipsoid}. Gauss developed the \href{https://en.wikipedia.org/wiki/Transverse_Mercator_projection\#Ellipsoidal_transverse_Mercator}{universal transverse Mercator projection} of the ellipsoidal shaped Earth (what he named \emph{conform projection}) for representing geodetical data in plane charts.

When the arc measurement was finished, Gauss began the enlargement of the triangulation to the west to get a survey of the whole \href{https://en.wikipedia.org/wiki/Kingdom_of_Hanover}{Kingdom of Hanover} with a Royal decree from 25 March 1828. The practical work was directed by three army officers, among them Lieutenant Joseph Gauss. The complete data evaluation laid in the hands of Gauss, who applied his mathematical inventions such as the \href{https://en.wikipedia.org/wiki/Method_of_least_squares}{method of least squares} and the \href{https://en.wikipedia.org/wiki/Gaussian_elimination}{elimination method} to it. The project was finished in 1844, and Gauss sent a final report of the project to the government; his method of projection was not edited until 1866.

In 1828, when studying differences in \href{https://en.wikipedia.org/wiki/Latitude}{latitude}, Gauss first defined a physical approximation for the \href{https://en.wikipedia.org/wiki/Figure_of_the_Earth}{figure of the Earth} as the surface everywhere perpendicular to the direction of gravity; later his doctoral student \href{https://en.wikipedia.org/wiki/Johann_Benedict_Listing}{Johann Benedict Listing} called this the \emph{\href{https://en.wikipedia.org/wiki/Geoid}{geoid}}.


\subsubsection{Magnetism and telegraphy}\label{Magnetism_and_telegraphy}


\paragraph{Geomagnetism}\label{Geomagnetism}

\begin{figure}
\centering
\href{https://en.wikipedia.org/wiki/File:Göttingen-Gauß-Weber-Monument.01.JPG}{\includegraphics[width=1.97917in,height=2.63542in]{250px-G_ttingen-Gau_-Weber-Monument.01.JPG}}
\caption{Gauss-Weber monument in Göttingen by Ferdinand Hartzer (1899)}\label{mwB8E}
\end{figure}

\begin{figure}
\centering
\href{https://en.wikipedia.org/wiki/File:A_magnetometer_used_by_Carl_Friedrich_Gauss,_from_Gerlach_und_F._Traumüller,_1899.png}{\includegraphics[width=2.60417in,height=1.67708in]{250px-A_magnetometer_used_by_Carl_Friedrich_Gauss_from_Gerlach_und_F._Traum_ller_1899.png}}
\caption{The Gauss--Weber magnetometer}\label{mwB8U}
\end{figure}

Gauss had been interested in magnetism since 1803. After \href{https://en.wikipedia.org/wiki/Alexander_von_Humboldt}{Alexander von Humboldt} visited Göttingen in 1826, both scientists began intensive research on \href{https://en.wikipedia.org/wiki/Geomagnetism}{geomagnetism}, partly independently, partly in productive cooperation. In 1828, Gauss was Humboldt\textquotesingle s guest during the conference of the \href{https://en.wikipedia.org/wiki/Society_of_German_Natural_Scientists_and_Physicians}{Society of German Natural Scientists and Physicians} in Berlin, where he got acquainted with the physicist \href{https://en.wikipedia.org/wiki/Wilhelm_Eduard_Weber}{Wilhelm Weber}.

When Weber got the chair for physics in Göttingen as successor of \href{https://en.wikipedia.org/wiki/Johann_Tobias_Mayer}{Johann Tobias Mayer} by Gauss\textquotesingle s recommendation in 1831, both of them started a fruitful collaboration, leading to a new knowledge of \href{https://en.wikipedia.org/wiki/Magnetism}{magnetism} with a representation for the unit of magnetism in terms of mass, charge, and time. They founded the \emph{Magnetic Association} (German: \emph{Magnetischer Verein}), an international working group of several observatories, which carried out measurements of \href{https://en.wikipedia.org/wiki/Earth's_magnetic_field}{Earth\textquotesingle s magnetic field} in many regions of the world using equivalent methods at arranged dates in the years 1836 to 1841.

In 1836, Humboldt suggested the establishment of a worldwide net of geomagnetic stations in the \href{https://en.wikipedia.org/wiki/British_Empire}{British dominions} with a letter to the \href{https://en.wikipedia.org/wiki/Prince_Augustus_Frederick,_Duke_of_Sussex}{Duke of Sussex}, then president of the Royal Society; he proposed that magnetic measures should be taken under standardized conditions using his methods. Together with other instigators, this led to a global program known as "\href{https://en.wikipedia.org/wiki/Edward_Sabine\#Magnetical_crusade}{Magnetical crusade}" under the direction of \href{https://en.wikipedia.org/wiki/Edward_Sabine}{Edward Sabine}. The dates, times, and intervals of observations were determined in advance, the \emph{Göttingen mean time} was used as the standard. 61 stations on all five continents participated in this global program. Gauss and Weber founded a series for publication of the results, six volumes were edited between 1837 and 1843. Weber\textquotesingle s departure to \href{https://en.wikipedia.org/wiki/Leipzig_University}{Leipzig} in 1843 as late effect of the \href{https://en.wikipedia.org/wiki/Göttingen_Seven}{Göttingen Seven affair} marked the end of Magnetic Association activity.

Following Humboldt\textquotesingle s example, Gauss ordered a magnetic \href{https://en.wikipedia.org/wiki/Observatory}{observatory} to be built in the garden of the observatory, but the scientists differed over instrumental equipment; Gauss preferred stationary instruments, which he thought to give more precise results, whereas Humboldt was accustomed to movable instruments. Gauss was interested in the temporal and spatial variation of magnetic \href{https://en.wikipedia.org/wiki/Magnetic_declination}{declination}, \href{https://en.wikipedia.org/wiki/Magnetic_dip}{inclination}, and intensity and differentiated, unlike Humboldt, between "horizontal" and "vertical" intensity. Together with Weber, he developed methods of measuring the components of the intensity of the magnetic field and constructed a suitable \href{https://en.wikipedia.org/wiki/Magnetometer}{magnetometer} to measure \emph{absolute values} of the strength of the Earth\textquotesingle s magnetic field, not more relative ones that depended on the apparatus. The precision of the magnetometer was about ten times higher than that of previous instruments. With this work, Gauss was the first to derive a non-mechanical quantity by basic mechanical quantities. He devised \href{https://en.wikipedia.org/wiki/Spherical_harmonics}{spherical harmonic} analysis as a technique to describe potential fields, and used it to show most of \href{https://en.wikipedia.org/wiki/Earth's_magnetic_field}{Earth\textquotesingle s magnetic field} was from internal sources.

Gauss carried out a \emph{General Theory of Terrestrial Magnetism} (1839), in what he believed to describe the nature of magnetic force; according to Felix Klein, this work is a presentation of observations by use of spherical harmonics rather than a physical theory. The theory predicted the existence of exactly two \href{https://en.wikipedia.org/wiki/Poles_of_astronomical_bodies\#Magnetic_poles}{magnetic poles} on the Earth, thus \href{https://en.wikipedia.org/wiki/Christopher_Hansteen}{Hansteen}\textquotesingle s idea of four magnetic poles became obsolete, and the data allowed to determine their location with rather good precision.

Gauss influenced the beginning of geophysics in Russia, when \href{https://en.wikipedia.org/wiki/Adolph_Theodor_Kupffer}{Adolph Theodor Kupffer}, one of his former students, founded a magnetic observatory in \href{https://en.wikipedia.org/wiki/St._Petersburg}{St. Petersburg}, following the example of the observatory in Göttingen, and similarly, \href{https://en.wikipedia.org/wiki/Ivan_Simonov}{Ivan Simonov} in \href{https://en.wikipedia.org/wiki/Kazan}{Kazan}.


\paragraph{Electromagnetism}\label{Electromagnetism}

\begin{figure}
\centering
\href{https://en.wikipedia.org/wiki/File:Gauss-Weber-Telegraf_Paulinerkirche_02.jpg}{\includegraphics[width=1.97917in,height=2.63542in]{250px-Gauss-Weber-Telegraf_Paulinerkirche_02.jpg}}
\caption{Town plan of Göttingen with course of the telegraphic connection}\label{mwCDY}
\end{figure}

The discoveries of \href{https://en.wikipedia.org/wiki/Hans_Christian_Ørsted}{Hans Christian Ørsted} on \href{https://en.wikipedia.org/wiki/Electromagnetism}{electromagnetism} and \href{https://en.wikipedia.org/wiki/Michael_Faraday}{Michael Faraday} on \href{https://en.wikipedia.org/wiki/Electromagnetic_induction}{electromagnetic induction} drew Gauss\textquotesingle s attention to these matters. Gauss and Weber found rules for branched \href{https://en.wikipedia.org/wiki/Electricity}{electric} circuits, which were later found independently and first published by \href{https://en.wikipedia.org/wiki/Gustav_Kirchhoff}{Gustav Kirchhoff} and named after him as \href{https://en.wikipedia.org/wiki/Kirchhoff's_circuit_laws}{Kirchhoff\textquotesingle s circuit laws}, and made inquiries into electromagnetism. They constructed the first \href{https://en.wikipedia.org/wiki/Electrical_telegraph}{electromechanical telegraph} in 1833, and Weber himself connected the observatory with the institute for physics in the town centre of Göttingen, but they made no further commercial use of this invention.

Gauss\textquotesingle s main theoretical interests in electromagnetism were reflected in his attempts to formulate quantitive laws governing electromagnetic induction. In notebooks from these years, he recorded several innovative formulations; he discovered the \href{https://en.wikipedia.org/wiki/Vector_potential}{vector potential} function, independently rediscovered by \href{https://en.wikipedia.org/wiki/Franz_Ernst_Neumann}{Franz Ernst Neumann} in 1845, and in January 1835 he wrote down an "induction law" equivalent to \href{https://en.wikipedia.org/wiki/Faraday's_law_of_induction}{Faraday\textquotesingle s law}, which stated that the \href{https://en.wikipedia.org/wiki/Electromotive_force}{electromotive force} at a given point in space is equal to the \href{https://en.wikipedia.org/wiki/Instantaneous_rate_of_change}{instantaneous rate of change} (with respect to time) of this function.

Gauss tried to find a unifying law for long-distance effects of \href{https://en.wikipedia.org/wiki/Electrostatics}{electrostatics}, \href{https://en.wikipedia.org/wiki/Electrodynamics}{electrodynamics}, electromagnetism, and \href{https://en.wikipedia.org/wiki/Electric_Induction}{induction}, comparable to Newton\textquotesingle s law of gravitation, but his attempt ended in a "tragic failure".


\subsubsection{Potential theory}\label{Potential_theory}

Since Isaac Newton had shown theoretically that the Earth and rotating stars assume non-spherical shapes, the problem of attraction of ellipsoids gained importance in mathematical astronomy. In his first publication on potential theory, the "Theoria attractionis..." (1813), Gauss provided a \href{https://en.wikipedia.org/wiki/Closed-form_expression}{closed-form expression} to the gravitational attraction of a homogeneous \href{https://en.wikipedia.org/wiki/Triaxial_ellipsoid}{triaxial ellipsoid} at every point in space. In contrast to previous research of \href{https://en.wikipedia.org/wiki/Colin_Maclaurin}{Maclaurin}, Laplace and Lagrange, Gauss\textquotesingle s new solution treated the attraction more directly in the form of an elliptic integral. In the process, he also proved and applied some special cases of the so-called \href{https://en.wikipedia.org/wiki/Divergence_theorem}{Gauss\textquotesingle s theorem} in \href{https://en.wikipedia.org/wiki/Vector_analysis}{vector analysis}.

In the \emph{General theorems concerning the attractive and repulsive forces acting in reciprocal proportions of quadratic distances} (1840) Gauss gave a basic theory of \href{https://en.wikipedia.org/wiki/Magnetic_vector_potential}{magnetic potential}, based on Lagrange, Laplace, and Poisson; it seems rather unlikely that he knew the previous works of \href{https://en.wikipedia.org/wiki/George_Green_(mathematician)}{George Green} on this subject. However, Gauss could never give any reasons for magnetism, nor a theory of magnetism similar to Newton\textquotesingle s work on gravitation, that enabled scientists to predict geomagnetic effects in the future.


\subsubsection{Optics}\label{Optics}

Gauss\textquotesingle s calculations enabled instrument maker \href{https://en.wikipedia.org/wiki/Johann_Georg_Repsold}{Johann Georg Repsold} in \href{https://en.wikipedia.org/wiki/Hamburg}{Hamburg} to construct a new \href{https://en.wikipedia.org/wiki/Achromatic_lens}{achromatic lens} system in 1810. A main problem, among other difficulties, was that the \href{https://en.wikipedia.org/wiki/Refractive_index}{refractive index} and \href{https://en.wikipedia.org/wiki/Dispersion_(optics)}{dispersion} of the glass used were not precisely known. In a short article from 1817 Gauss dealt with the problem of removal of \href{https://en.wikipedia.org/wiki/Chromatic_aberration}{chromatic aberration} in \href{https://en.wikipedia.org/wiki/Gauss_lens}{double lenses}, and computed adjustments of the shape and coefficients of refraction required to minimize it. His work was noted by the optician \href{https://en.wikipedia.org/wiki/Carl_August_von_Steinheil}{Carl August von Steinheil}, who in 1860 introduced the achromatic \href{https://en.wikipedia.org/wiki/Achromatic_lens}{Steinheil doublet}, partly based on Gauss\textquotesingle s calculations. Many results in \href{https://en.wikipedia.org/wiki/Geometrical_optics}{geometrical optics} are scattered in Gauss\textquotesingle s correspondences and hand notes.

In the \emph{Dioptrical Investigations} (1840), Gauss gave the first systematic analysis of the formation of images under a \href{https://en.wikipedia.org/wiki/Paraxial_approximation}{paraxial approximation} (\href{https://en.wikipedia.org/wiki/Gaussian_optics}{Gaussian optics}). He characterized optical systems under a paraxial approximation only by its \href{https://en.wikipedia.org/wiki/Cardinal_point_(optics)}{cardinal points}, and he derived the Gaussian \href{https://en.wikipedia.org/wiki/Lens}{lens} formula, applicable without restrictions in respect to the thickness of the lenses.


\subsubsection{Mechanics}\label{Mechanics}

Gauss\textquotesingle s first work in mechanics concerned the \href{https://en.wikipedia.org/wiki/Earth's_rotation}{earth\textquotesingle s rotation}. When his university friend \href{https://en.wikipedia.org/wiki/Johann_Benzenberg}{Benzenberg} carried out experiments to determine the deviation of falling masses from the perpendicular in 1802, what today is known as the \href{https://en.wikipedia.org/wiki/Coriolis_force}{Coriolis force}, he asked Gauss for a theory-based calculation of the values for comparison with the experimental ones. Gauss elaborated a system of fundamental equations for the motion, and the results corresponded sufficiently with Benzenberg\textquotesingle s data, who added Gauss\textquotesingle s considerations as an appendix to his book on falling experiments.

After \href{https://en.wikipedia.org/wiki/Léon_Foucault}{Foucault} had demonstrated the earth\textquotesingle s rotation by his \href{https://en.wikipedia.org/wiki/Foucault_pendulum}{pendulum} experiment in public in 1851, Gerling questioned Gauss for further explanations. This instigated Gauss to design a new apparatus for demonstration with a much shorter length of pendulum than Foucault\textquotesingle s one. The oscillations were observed with a reading telescope, with a vertical scale and a mirror fastened at the pendulum. It is described in the Gauss--Gerling correspondence and Weber made some experiments with this apparatus in 1853, but no data were published.

\href{https://en.wikipedia.org/wiki/Gauss's_principle_of_least_constraint}{Gauss\textquotesingle s principle of least constraint} of 1829 was established as a general concept to overcome the division of mechanics into statics and dynamics, combining \href{https://en.wikipedia.org/wiki/D'Alembert's_principle}{D\textquotesingle Alembert\textquotesingle s principle} with \href{https://en.wikipedia.org/wiki/Joseph-Louis_Lagrange}{Lagrange}\textquotesingle s \href{https://en.wikipedia.org/wiki/Principle_of_virtual_work}{principle of virtual work}, and showing analogies to the method of \href{https://en.wikipedia.org/wiki/Least_squares}{least squares}.


\subsubsection{Metrology}\label{Metrology}

In 1828, Gauss was appointed as head of the board for weights and measures of the Kingdom of Hanover. He created \href{https://en.wikipedia.org/wiki/Standard_(metrology)}{standards} for length and measure. Gauss himself took care of the time-consuming measures and gave detailed orders for the mechanical construction. In the correspondence with Schumacher, who was also working on this matter, he described new ideas for high-precision scales. He submitted the final reports on the Hanoverian \href{https://en.wikipedia.org/wiki/Foot_(unit)}{foot} and \href{https://en.wikipedia.org/wiki/Pound_(mass)}{pound} to the government in 1841. This work achieved international importance due to an 1836 law that connected the Hanoverian measures with the English ones.


\subsection{Honours and awards}\label{Honours_and_awards}

\begin{figure}
\centering
\href{https://en.wikipedia.org/wiki/File:Carl_Friedrich_Gauß,_Copley-Medaille,_1838.jpg}{\includegraphics[width=1.97917in,height=3.69792in]{250px-Carl_Friedrich_Gau_Copley-Medaille_1838.jpg}}
\caption{\href{https://en.wikipedia.org/wiki/Copley_Medal}{Copley Medal} for Gauss (1838)}\label{mwCO4}
\end{figure}

Gauss first became member of a scientific society, the \href{https://en.wikipedia.org/wiki/Russian_Academy_of_Sciences}{Russian Academy of Sciences}, in 1802. Further memberships (corresponding, foreign or full) were awarded by the \href{https://en.wikipedia.org/wiki/Göttingen_Academy_of_Sciences_and_Humanities}{Academy of Sciences} in Göttingen (1802/ 1807), the \href{https://en.wikipedia.org/wiki/French_Academy_of_Sciences}{French Academy of Sciences} (1804/ 1820), the \href{https://en.wikipedia.org/wiki/Royal_Society}{Royal Society} of London (1804), the \href{https://en.wikipedia.org/wiki/Prussian_Academy_of_Sciences}{Royal Prussian Academy} in Berlin (1810), the \href{https://en.wikipedia.org/wiki/Accademia_nazionale_delle_scienze}{National Academy of Science} in Verona (1810), the \href{https://en.wikipedia.org/wiki/Royal_Society_of_Edinburgh}{Royal Society of Edinburgh} (1820), the \href{https://en.wikipedia.org/wiki/Bavarian_Academy_of_Sciences_and_Humanities}{Bavarian Academy of Sciences} of Munich (1820), the \href{https://en.wikipedia.org/wiki/Royal_Danish_Academy_of_Sciences_and_Letters}{Royal Danish Academy} in Copenhagen (1821), the \href{https://en.wikipedia.org/wiki/Royal_Astronomical_Society}{Royal Astronomical Society} in London (1821), the \href{https://en.wikipedia.org/wiki/Royal_Swedish_Academy_of_Sciences}{Royal Swedish Academy of Sciences} (1821), the \href{https://en.wikipedia.org/wiki/American_Academy_of_Arts_and_Sciences}{American Academy of Arts and Sciences} in Boston (1822), the \href{https://en.wikipedia.org/wiki/Royal_Bohemian_Society_of_Sciences}{Royal Bohemian Society of Sciences} in Prague (1833), the \href{https://en.wikipedia.org/wiki/Royal_Academy_of_Science,_Letters_and_Fine_Arts_of_Belgium}{Royal Academy of Science, Letters and Fine Arts of Belgium} (1841/1845), the \href{https://en.wikipedia.org/wiki/Royal_Society_of_Sciences_in_Uppsala}{Royal Society of Sciences in Uppsala} (1843), the \href{https://en.wikipedia.org/wiki/Royal_Irish_Academy}{Royal Irish Academy} in Dublin (1843), the \href{https://en.wikipedia.org/wiki/Royal_Netherlands_Academy_of_Arts_and_Sciences}{Royal Institute of the Netherlands} (1845/ 1851), the \href{https://en.wikipedia.org/wiki/Spanish_Royal_Academy_of_Sciences}{Spanish Royal Academy of Sciences} in Madrid (1850), the \href{https://en.wikipedia.org/wiki/Russian_Geographical_Society}{Russian Geographical Society} (1851), the \href{https://en.wikipedia.org/wiki/Austrian_Academy_of_Sciences}{Imperial Academy of Sciences} in Vienna (1848), the \href{https://en.wikipedia.org/wiki/American_Philosophical_Society}{American Philosophical Society} (1853), the \href{https://en.wikipedia.org/wiki/Cambridge_Philosophical_Society}{Cambridge Philosophical Society}, and the \href{https://en.wikipedia.org/wiki/Koninklijke_Hollandsche_Maatschappij_der_Wetenschappen}{Royal Hollandish Society of Sciences} in Haarlem.

Both the \href{https://en.wikipedia.org/wiki/University_of_Kazan}{University of Kazan} and the Philosophy Faculty of the \href{https://en.wikipedia.org/wiki/Charles_University}{University of Prague} appointed him honorary member in 1848.

Gauss received the \href{https://en.wikipedia.org/wiki/Lalande_Prize}{Lalande Prize} from the French Academy of Science in 1809 for the theory of planets and the means of determining their orbits from only three observations, the Danish Academy of Science prize in 1823 for his memoir on conformal projection, and the \href{https://en.wikipedia.org/wiki/Copley_Medal}{Copley Medal} from the Royal Society in 1838 for "his inventions and mathematical researches in magnetism".

Gauss was appointed Knight of the French \href{https://en.wikipedia.org/wiki/Legion_of_Honour}{Legion of Honour} in 1837, and became one of the first members of the Prussian \href{https://en.wikipedia.org/wiki/Order_Pour_le_Merite\#Civil_class}{Order Pour le Merite} (Civil class) when it was established in 1842. He received the \href{https://en.wikipedia.org/wiki/Order_of_the_Crown_of_Westphalia}{Order of the Crown of Westphalia} (1810), the Danish \href{https://en.wikipedia.org/wiki/Order_of_the_Dannebrog}{Order of the Dannebrog} (1817), the Hanoverian \href{https://en.wikipedia.org/wiki/Royal_Guelphic_Order}{Royal Guelphic Order} (1815), the Swedish \href{https://en.wikipedia.org/wiki/Order_of_the_Polar_Star}{Order of the Polar Star} (1844), the \href{https://en.wikipedia.org/wiki/Order_of_Henry_the_Lion}{Order of Henry the Lion} (1849), and the \href{https://en.wikipedia.org/wiki/Bavarian_Maximilian_Order_for_Science_and_Art}{Bavarian Maximilian Order for Science and Art} (1853).

The Kings of Hanover appointed him the honorary titles "\href{https://en.wikipedia.org/wiki/Hofrath}{Hofrath}" (1816) and "Geheimer Hofrath" (1845). In 1949, on the occasion of his golden doctor degree jubilee, he received \href{https://en.wikipedia.org/wiki/Honorary_citizenship}{honorary citizenship} of both Brunswick and Göttingen. Soon after his death a medal was issued by order of King \href{https://en.wikipedia.org/wiki/George_V_of_Hanover}{George V of Hanover} with the back inscription dedicated "to the Prince of Mathematicians".

The "Gauss-Gesellschaft Göttingen" ("Göttingen Gauss Society") was founded in 1964 for research on the life and work of Carl Friedrich Gauss and related persons. It publishes the \emph{Mitteilungen der Gauss-Gesellschaft} (\emph{Communications of the Gauss Society}).


\subsection{Names and commemorations}\label{Names_and_commemorations}

\begin{itemize}
\tightlist
\item
  {\href{https://en.wikipedia.org/wiki/List_of_things_named_after_Carl_Friedrich_Gauss}{List of things named after Carl Friedrich Gauss}}
\end{itemize}


\subsection{Selected writings}\label{Selected_writings}


\subsubsection{Mathematics and astronomy}\label{Mathematics_and_astronomy}

\begin{figure}
\centering
\href{https://en.wikipedia.org/wiki/File:Gauß-Denkmal_-_Braunschweig,_Germany_-_DSC04436.JPG}{\includegraphics[width=1.97917in,height=2.9375in]{250px-Gau_-Denkmal_-_Braunschweig_Germany_-_DSC04436.JPG}}
\caption{Statue of Gauss in \href{https://en.wikipedia.org/wiki/Braunschweig}{Brunswick} (1880), made by \href{https://en.wikipedia.org/wiki/Hermann_Heinrich_Howaldt}{Hermann Heinrich Howaldt}, designed by \href{https://en.wikipedia.org/wiki/Fritz_Schaper}{Fritz Schaper}}\label{mwCdM}
\end{figure}

\begin{itemize}
\tightlist
\item
  {1799: Demonstratio nova theorematis omnem functionem algebraicam rationalem integram unius variabilis in factores reales primi vel secundi gradus resolvi posse {{[}}\emph{New proof of the theorem that every integral algebraic function of one variable can be resolved into real factors of the first or second degree}{{]}}. Helmstedt: C. G. Fleckeisen.{} (Doctoral thesis on the fundamental theorem of algebra, University of Helmstedt) \href{https://edoc.hu-berlin.de/handle/18452/732}{Original book}}
\item
  {1816: "Demonstratio nova altera theorematis omnem functionem algebraicam rationalem integram unius variabilis in factores reales primi vel secundi gradus resolvi posse". \emph{Commentationes Societatis Regiae Scientiarum Gottingensis Recentiores. Comm. Class. Math}. \textbf{3}: {107--}134.{} Original}
\item
  {1816: "Theorematis de resolubilitate functionum algebraicarum integrarum in factores reales demonstratio tertia". \emph{Commentationes Societatis Regiae Scientiarum Gottingensis Recentiores. Comm. Class. Math}. \textbf{3}: {135--}142.{} Original}
\item
  {1850: "Beiträge zur Theorie der algebraischen Gleichungen". \emph{Abhandlungen der Königlichen Gesellschaft der Wissenschaften zu Göttingen}. \textbf{4}: {34--}35.{} Original (Lecture from 1849)}

  \begin{itemize}
  \tightlist
  \item
    {\href{https://archive.org/details/bub_gb_LoW4AAAAIAAJ/mode/2up}{\emph{Die vier Gauss\textquotesingle schen Beweise für die Zerlegung ganzer algebraischer Funktionen in reelle Faktoren ersten und zweiten Grades. (1799--1849)}} {{[}}\emph{The four Gaussian proofs of the fundamental theorem of algebra}{{]}}. Translated by \href{https://en.wikipedia.org/wiki/Eugen_Netto}{Netto}. Leipzig: Wilhelm Engelmann. 1890.{} (German)}
  \end{itemize}
\item
  {1800: "Berechnung des Osterfestes" {{[}}Calculation of Easter{{]}}. \emph{Monatliche Correspondenz zur Beförderung der Erd- und Himmelskunde} (in German). \textbf{2}: {121--}130.{} \href{https://zs.thulb.uni-jena.de/receive/jportal_jpvolume_00201970}{Original}}
\item
  {1801: Disquisitiones Arithmeticae. Leipzig: Gerh. Fleischer jun.{}}

  \begin{itemize}
  \tightlist
  \item
    {Gauss, Carl Friedrich (1986). \href{https://link.springer.com/book/10.1007/978-1-4939-7560-0\#about-this-book}{\emph{Disquisitiones Arithmeticae \& other papers on number theory}}. Translated by Clarke, Arthur A. (2nd, corrected{~}ed.). New York: \href{https://en.wikipedia.org/wiki/Springer_Science+Business_Media}{Springer}. \href{https://en.wikipedia.org/wiki/Doi_(identifier)}{doi}:\href{https://doi.org/10.1007\%2F978-1-4939-7560-0}{10.1007/978-1-4939-7560-0}. \href{https://en.wikipedia.org/wiki/ISBN_(identifier)}{ISBN}{~}\href{https://en.wikipedia.org/wiki/Special:BookSources/978-0-387-96254-2}{978-0-387-96254-2}.{} (translated from the \href{https://www.science.org/doi/10.1126/science.154.3749.642.b}{second German edition, Göttingen 1860})}
  \end{itemize}
\item
  {1802: "Berechnung des jüdischen Osterfestes" {{[}}Calculation of Jewish Easter{{]}}. \emph{Monatliche Correspondenz zur Beförderung der Erd- und Himmelskunde} (in German). \textbf{5}: {435--}437.{} \href{https://zs.thulb.uni-jena.de/rsc/viewer/jportal_derivate_00237503/Monatlich_Correspondenz_130168688_5_1802_0444\%20.tif}{Original}}
\item
  {1804: "Über die Grenzen der geocentrischen Oerter der Planeten" {{[}}On the limits of the geocentric places of the planets{{]}}. \emph{Monatliche Correspondenz zur Beförderung der Erd- und Himmelskunde} (in German). \textbf{10}: {171--}193.{} \href{https://zs.thulb.uni-jena.de/rsc/viewer/jportal_derivate_00240854/Monatliche_Correspondenz_130168688_10_0175.tif}{Original} (on the Zodiacus)}
\item
  {1808: "Theorematis arithmetici demonstratio nova". \emph{Commentationes Societatis Regiae Scientiarum Gottingensis. Comm. Math}. \textbf{16}: {69--}74.{} Original (Introduces \href{https://en.wikipedia.org/wiki/Gauss's_lemma_(number_theory)}{Gauss\textquotesingle s lemma}, uses it in the third proof of quadratic reciprocity)}
\item
  {1808: Methodus peculiaris elevationem poli determinandi (in Latin). Göttingen.{}}
\item
  {1809: Theoria motus corporum coelestium in sectionibus conicis solem ambientium (in Latin). Hamburg: Friedrich Perthes \& Johann Heinrich Besser.{} \href{https://gutenberg.beic.it/webclient/DeliveryManager?pid=12217180}{Original book}}

  \begin{itemize}
  \tightlist
  \item
    {\href{https://archive.org/details/theoryofmotionof00gausrich/}{\emph{Theory of the Motion of Heavenly Bodies Moving about the Sun in Conic Sections}}. Translated by \href{https://en.wikipedia.org/wiki/Charles_Henry_Davis}{Davis, Charles Henry}. Little, Brown \& Co. 1857.{}}
  \item
    {\emph{Theory of the motion of the celestial bodies moving around the Sun in conic sections. Reprint of the 1809 original. (Theoria motus corporum coelestium in sectionibus conicis solem ambientium.) (Latin)}. Cambridge Library Collection - Mathematics. \href{https://en.wikipedia.org/wiki/Cambridge_University_Press}{Cambridge University Press}. 2011. \href{https://en.wikipedia.org/wiki/ISBN_(identifier)}{ISBN}{~}\href{https://en.wikipedia.org/wiki/Special:BookSources/978-1-108-14311-0}{978-1-108-14311-0}. \href{https://en.wikipedia.org/wiki/Zbl_(identifier)}{Zbl}{~}\href{https://zbmath.org/?format=complete&q=an:1234.01016}{1234.01016}.{}}
  \end{itemize}
\item
  {1811: "Disquisitio de elementis ellipticis Palladis ex oppositionibus annorum 1803, 1804, 1805, 1806, 1807, 1808, 1809". \emph{Commentationes Societatis Regiae Scientiarum Gottingensis Recentiores. Comm. Math}. \textbf{1}: {1--}26.{} Original (from 1810) (Orbit of \href{https://en.wikipedia.org/wiki/2_Pallas}{Pallas})}
\item
  {1811: "Summatio quarundam serierum singularium". \emph{Commentationes Societatis Regiae Scientiarum Gottingensis Recentiores. Comm. Class. Math}. \textbf{1}: {1--}40.{} Original (from 1808) (Determination of the sign of the \href{https://en.wikipedia.org/wiki/Quadratic_Gauss_sum}{quadratic Gauss sum}, uses this to give the fourth proof of quadratic reciprocity)}
\item
  {1813: "Disquisitiones generales circa seriem infinitam 1 + α β γ .1 + etc. \{\textbackslash displaystyle 1+\{\textbackslash frac \{\textbackslash alpha \textbackslash beta \}\{\textbackslash gamma .1\}\}+\{\textbackslash mbox\{etc.\}\}\} ". \emph{Commentationes Societatis Regiae Scientiarum Gottingensis Recentiores. Comm. Class. Math}. \textbf{2}: {1--}42.{} Original (from 1812, contains the \href{https://en.wikipedia.org/wiki/Gauss's_continued_fraction}{Gauss\textquotesingle s continued fraction})}
\item
  {1816: "Methodus nova integralium valores per approximationem inveniendi". \emph{Commentationes Societatis Regiae Scientiarum Gottingensis Recentiores. Comm. Class. Math}. \textbf{3}: {39--}76.{} Original (from 1814)}
\item
  {1818: "Theorematis fundamentalis in doctrina de residuis quadraticis demonstrationes et ampliationes novae". \emph{Commentationes Societatis Regiae Scientiarum Gottingensis Recentiores. Comm. Class. Math}. \textbf{4}: {3--}20.{} Original (from 1817) (Fifth and sixth proofs of quadratic reciprocity)}
\item
  {1818: "Determinatio attractionis, quam in punctum positionis datae exerceret planeta, si eius massa per totamorbitam, ratione temporis, quo singulae partes describuntur, uniformiter esset dispertita". \emph{Commentationes Societatis Regiae Scientiarum Gottingensis Recentiores. Comm. Class. Math}. \textbf{4}: {21--}48.{} Original (Only reference to the -- mostly unpublished -- work on the algorithm of the \href{https://en.wikipedia.org/wiki/Arithmetic-geometric_mean}{arithmetic-geometric mean}.)}
\item
  {1823: "Theoria combinationis observationum erroribus minimis obnoxiae. Pars Prior". \emph{Commentationes Societatis Regiae Scientiarum Gottingensis Recentiores. Comm. Class. Math}. \textbf{5}: {33--}62.{} Original (from 1821)}
\item
  {1823: "Theoria combinationis observationum erroribus minimis obnoxiae. Pars Posterior". \emph{Commentationes Societatis Regiae Scientiarum Gottingensis Recentiores. Comm. Class. Math}. \textbf{5}: {63--}90.{} Original}
\item
  {1825: "Allgemeine Auflösung der Aufgabe die Theile einer gegebnen Fläche auf einer andern gegebnen Fläche so abzubilden dass die Abbildung dem Abgebildeten in den kleinsten Theilen ähnlich wird". \emph{Astronomische Abhandlungen}. \textbf{3}. Altona.{} (Prize winning essay from 1822 on conformal mapping)}
\item
  {1828: Bestimmung des Breitenunterschiedes zwischen den Sternwarten von Göttingen und Altona durch Beobachtungen am Ramsdenschen Zenithsector {{[}}\emph{Determination of the Difference in Latitude between the Observatories of Göttingen and Altona by Observations with Ramsden\textquotesingle s Zenith sector}{{]}} (in German). Göttingen: Vandenhoeck und Ruprecht. 1828.{} \href{https://books.google.com/books?id=tIg_AAAAcAAJ&pg=PA04}{Original book}}
\item
  {1828: Gauss, Carl Friedrich (1828). "Supplementum theoriae combinationis observationum erroribus minimis obnoxiae". \emph{Commentationes Societatis Regiae Scientiarum Gottingensis Recentiores. Comm. Class. Math}. \textbf{6}: {57--}98. \href{https://en.wikipedia.org/wiki/Bibcode_(identifier)}{Bibcode}:\href{https://ui.adsabs.harvard.edu/abs/1828stco.book.....G}{1828stco.book.....G}.{} (from 1826)}

  \begin{itemize}
  \tightlist
  \item
    {Gauss, Carl Friedrich; Stewart, G. W. (1995). \href{https://epubs.siam.org/doi/book/10.1137/1.9781611971248}{\emph{Theory of the Combination of Observations Least Subject to Errors. Part One, Part Two, Supplement (Classics in Applied Mathematics)}}. Translated by G. W. Stewart. Philadelphia: Society for Industrial and Applied Mathematics. \href{https://en.wikipedia.org/wiki/Doi_(identifier)}{doi}:\href{https://doi.org/10.1137\%2F1.9781611971248}{10.1137/1.9781611971248}. \href{https://en.wikipedia.org/wiki/ISBN_(identifier)}{ISBN}{~}\href{https://en.wikipedia.org/wiki/Special:BookSources/978-0-89871-347-3}{978-0-89871-347-3}.{} (Three essays concerning the calculation of probabilities as the basis of the Gaussian law of error propagation)}
  \end{itemize}
\item
  {1828: "Disquisitiones generales circa superficies curvas". \emph{Commentationes Societatis Regiae Scientiarum Gottingensis Recentiores. Comm. Class. Math}. \textbf{6}: {99--}146.{} Original (from 1827)}

  \begin{itemize}
  \tightlist
  \item
    {\href{https://www.gutenberg.org/files/36856/36856-pdf.pdf}{\emph{General Investigations of Curved Surfaces}} {(PDF)}. Translated by J. C. Morehead and A. M. Hiltebeitel. The Princeton University Library. 1902.{}}
  \end{itemize}
\item
  {1828: "Theoria residuorum biquadraticorum, Commentatio prima". \emph{Commentationes Societatis Regiae Scientiarum Gottingensis Recentiores. Comm. Class. Math}. \textbf{6}: {27--}56.{} Original (from 1825)}
\item
  {1832: "Theoria residuorum biquadraticorum, Commentatio secunda". \emph{Commentationes Societatis Regiae Scientiarum Gottingensis Recentiores. Comm. Class. Math}. \textbf{7}: {89--}148.{} Original (from 1831) (Introduces the \href{https://en.wikipedia.org/wiki/Gaussian_integers}{Gaussian integers}, states (without proof) the law of \href{https://en.wikipedia.org/wiki/Biquadratic_reciprocity}{biquadratic reciprocity}, proves the supplementary law for 1 + \emph{i})}
\item
  {1845: "Untersuchungen über Gegenstände der Höheren Geodäsie. Erste Abhandlung". \emph{Abhandlungen der Königlichen Gesellschaft der Wissenschaften in Göttingen}. Zweiter Band, von den Jahren 1842--1844: {3--}46.{} \href{https://gdz.sub.uni-goettingen.de/id/PPN250442582_0002}{Original} (from 1843)}
\item
  {1847: "Untersuchungen über Gegenstände der Höheren Geodäsie. Zweite Abhandlung". \emph{Abhandlungen der Königlichen Gesellschaft der Wissenschaften in Göttingen}. Dritter Band, von den Jahren 1845--1847: {3--}44.{} \href{https://gdz.sub.uni-goettingen.de/id/PPN250442582_0003}{Original} (from 1846)}
\item
  {1848: Gauss (1848). "Schreiben des Herrn Geheimen Hofrathes Gauss an den Herausgeber" {{[}}Letter of Mr. Secret Councillor of the Court Gauss to the editor{{]}}. \emph{\href{https://en.wikipedia.org/wiki/Astronomische_Nachrichten}{Astronomische Nachrichten}} (in German). \textbf{27}: {1--}3. \href{https://en.wikipedia.org/wiki/Bibcode_(identifier)}{Bibcode}:\href{https://ui.adsabs.harvard.edu/abs/1848AN.....27....1G}{1848AN.....27....1G}. \href{https://en.wikipedia.org/wiki/Doi_(identifier)}{doi}:\href{https://doi.org/10.1002\%2Fasna.18480270102}{10.1002/asna.18480270102}.{} Original}
\item
  {1903: Wissenschaftliches Tagebuch (Klein, Felix, ed. (1903). "Gauß\textquotesingle{} wissenschaftliches Tagebuch 1796--1814". \emph{\href{https://en.wikipedia.org/wiki/Mathematische_Annalen}{Mathematische Annalen}} (in Latin and German). \textbf{57}: {1--}34. \href{https://en.wikipedia.org/wiki/Doi_(identifier)}{doi}:\href{https://doi.org/10.1007\%2FBF01449013}{10.1007/BF01449013}. \href{https://en.wikipedia.org/wiki/S2CID_(identifier)}{S2CID}{~}\href{https://api.semanticscholar.org/CorpusID:119641638}{119641638}.{}) \href{https://gdz.sub.uni-goettingen.de/id/DE-611-HS-3382323}{Original book} (from 1847, on the Zodiacus)}

  \begin{itemize}
  \tightlist
  \item
    {\href{https://en.wikipedia.org/wiki/Jeremy_Gray_(mathematician)}{Jeremy Gray} (1984). "A commentary on Gauss\textquotesingle s mathematical diary, 1796--1814". \emph{Expositiones Mathematicae}. \textbf{2}: {97--}130.{}}
  \end{itemize}
\end{itemize}


\subsubsection{Physics}\label{Physics}

\begin{itemize}
\tightlist
\item
  {1804: Fundamentalgleichungen für die Bewegung schwerer Körper auf der Erde ( in original book: \href{https://en.wikipedia.org/wiki/Johann_Benzenberg}{Benzenberg, Johann Friedrich}. \emph{Versuche über das Gesetz des Falls, über den Widerstand der Luft und über die Umdrehung der Erde} {{[}}\emph{Experiments on the Law of falling Bodies, on the Resistance of Air, and of the Rotation of the Earth}{{]}}. Dortmund: Gebrüder Mallinckrodt. pp.{~}{363--}371.{} \href{https://www.digitale-sammlungen.de/de/view/bsb10060427?page=392,393}{Original})}
\item
  {1813: "Theoria attractionis corporum sphaeroidicorum ellipticorum homogeneorum methodo nova tractata". \emph{Commentationes Societatis Regiae Scientiarum Gottingensis Recentiores. Comm. Class. Math}. \textbf{2}: {1--}24.{} Original (contains \href{https://en.wikipedia.org/wiki/Gauss's_theorem}{Gauss\textquotesingle s theorem} of vector analysis)}
\item
  {1817: \href{https://babel.hathitrust.org/cgi/pt?id=hvd.hxigq8&seq=349}{"Ueber die achromatischen Doppelobjective besonders in Rücksicht der vollkommnern Aufhebung der Farbenzerstreuung"} {{[}}On achromatic double lenses with special regard to a more complete dispersion of colours{{]}}. \emph{Zeitschrift für Astronomie und verwandte Wissenschaften} (in German). \textbf{IV}: {345--}351.{}}
\item
  {1829: \href{https://zenodo.org/record/1448816}{"Über ein neues allgemeines Grundgesetz der Mechanik"} {{[}}On a new General Fundamental Law of Mechanics{{]}}. \emph{Journal für die reine und angewandte Mathematik}. \textbf{1829} (4): {232--}235. 1829. \href{https://en.wikipedia.org/wiki/Doi_(identifier)}{doi}:\href{https://doi.org/10.1515\%2Fcrll.1829.4.232}{10.1515/crll.1829.4.232}. \href{https://en.wikipedia.org/wiki/S2CID_(identifier)}{S2CID}{~}\href{https://api.semanticscholar.org/CorpusID:199545985}{199545985}.{}}
\item
  {1830: "Principia generalia theoriae figurae fluidorum in statu aequilibrii". \emph{Commentationes Societatis Regiae Scientiarum Gottingensis Recentiores}. \textbf{7}: {39--}88.{} Original (from 1829)}
\item
  {1841: "Intensitas vis magneticae terrestris ad mensuram absolutam revocata". \emph{Commentationes Societatis Regiae Scientiarum Gottingensis Recentiores}. \textbf{8}: {3--}44.{} Original (from 1832)}

  \begin{itemize}
  \tightlist
  \item
    {\href{http://21sci-tech.com/translations/gaussMagnetic.pdf}{The Intensity of the Earth\textquotesingle s Magnetic Force Reduced to Absolute Measurement.} Translated by Susan P. Johnson.}
  \end{itemize}
\item
  {1836: Erdmagnetismus und Magnetometer (Original book: H.C. Schumacher (ed.). \href{https://www.digitale-sammlungen.de/de/view/bsb10538569?page=21}{\emph{Jahrbuch für 1836}} (in German). Vol.{~}1836. Tübingen: J.G.Cotta\textquotesingle sche Buchhandlung. pp.{~}{1--}47.{})}
\item
  {1840: Allgemeine Lehrsätze in Beziehung auf die im verkehrten Verhältnis des Quadrats der Entfernung wirkenden Anziehungs- und Abstoßungskräfte (Original book: \href{https://www.deutschestextarchiv.de/book/show/gauss_lehrsaetze_1840}{\emph{Allgemeine Lehrsätze in Beziehung auf die im verkehrten Verhältnis des Quadrats der Entfernung wirkenden Anziehungs- und Abstoßungskräfte}} {{[}}\emph{General Theorems concerning the attractive and repulsive Forces acting in reciprocal Proportions of quadratic Distances}{{]}} (in German). Leipzig: Weidmannsche Buchhandlung. 1840.{}}
\item
  {1843: "Dioptrische Untersuchungen" {{[}}Dioptrical Investigations{{]}}. \emph{Abhandlungen der Königlichen Gesellschaft der Wissenschaften in Göttingen} (in German). Erster Band: {1--}34.{} Original (from 1840)}
\end{itemize}


\paragraph{Together with Wilhelm Weber}\label{Together_with_Wilhelm_Weber}

\begin{itemize}
\tightlist
\item
  {1837--1839: Weber, Wilhelm Eduard; Gauss, Carl Friedrich. \href{https://babel.hathitrust.org/cgi/pt?id=njp.32101058086446&view=1up&seq=9}{\emph{Resultate aus den Beobachtungen des magnetischen Vereins im Jahre 1836--1838}} (in German). Göttingen: Dieterichsche Buchhandlung. pp.{~}6 v.{}}
\item
  {1840--1843: Weber, Wilhelm Eduard; Gauss, Carl Friedrich. \href{https://babel.hathitrust.org/cgi/pt?id=njp.32101058086396&view=1up&seq=199}{\emph{Resultate aus den Beobachtungen des magnetischen Vereins im Jahre 1839--1841}} (in German). Leipzig: Weidmannsche Verlagsbuchhandlung. pp.{~}6 v.{}}
\item
  {1840: Weber, Wilhelm Eduard; Gauss, Carl Friedrich. \href{https://babel.hathitrust.org/cgi/pt?id=njp.32101058086388&view=1up&seq=5}{\emph{Atlas des Erdmagnetismus nach den Elementen der Theorie entworfen. Supplement zu den Resultaten aus den Beobachtungen des magnetischen Vereins}} (in German). Leipzig: Weidmannsche Verlagsbuchhandlung. pp.{~}6 v.{}}
\end{itemize}


\subsubsection{Collected works}\label{Collected_works}

\begin{itemize}
\tightlist
\item
  {Königlich Preußische Akademie der Wissenschaften, ed. (1863--1933). Carl Friedrich Gauss. Werke (in Latin and German). Vol.{~}{1--}12. Göttingen: (diverse publishers).{} (includes unpublished literary estate)}
\end{itemize}


\subsubsection{Correspondence}\label{Correspondence}

\begin{itemize}
\tightlist
\item
  {Königlich Preußische Akademie der Wissenschaften, ed. (1880). \href{https://archive.org/details/briefwechselzwi00berlgoog/page/n5/mode/2up}{\emph{Briefwechsel zwischen Gauss und Bessel}} (in German). Leipzig: Wilhelm Engelmann.{} (letters from December 1804 to August 1844)}
\item
  {Schoenberg, Erich; Perlick, Alfons (1955). \emph{Unbekannte Briefe von C. F. Gauß und Fr. W. Bessel}. Abhandlungen der Bayerischen Akademie der Wissenschaften, Math.-nat. Klasse, Neue Folge, No. 71 (in German). Munich: Verlag der Bayerischen Akademie der Wissenschaften. pp.{~}{5--}21.{} (letters to \href{https://en.wikipedia.org/wiki/Palm_Heinrich_Ludwig_von_Boguslawski}{Boguslawski} from February 1835 to January 1848)}
\item
  {Schwemin, Friedhelm, ed. (2014). \emph{Der Briefwechsel zwischen Carl Friedrich Gauß und \href{https://en.wikipedia.org/wiki/Johann_Elert_Bode}{Johann Elert Bode}}. Acta Historica Astronomica (in German). Vol.{~}53. Leipzig: Akademische Verlaganstalt. \href{https://en.wikipedia.org/wiki/ISBN_(identifier)}{ISBN}{~}\href{https://en.wikipedia.org/wiki/Special:BookSources/978-3-944913-43-8}{978-3-944913-43-8}.{} (letters from February 1802 to October 1826)}
\item
  {Franz Schmidt, Paul Stäckel, ed. (1899). \href{https://archive.org/details/briefwechselzwi00gausgoog/page/n4/mode/2up}{\emph{Briefwechsel zwischen Carl Friedrich Gauss und Wolfgang Bolyai}} (in German). Leipzig: B.G. Teubner.{} (letters from September 1797 to February 1853; added letters of other correspondents)}
\item
  {Axel Wittmann, ed. (2018). \emph{Obgleich und indeßen. Der Briefwechsel zwischen Carl Friedrich Gauss und \href{https://en.wikipedia.org/wiki/Johann_Franz_Encke}{Johann Franz Encke}} (in German). Remagen: Verlag Kessel. \href{https://en.wikipedia.org/wiki/ISBN_(identifier)}{ISBN}{~}\href{https://en.wikipedia.org/wiki/Special:BookSources/978-3945941379}{978-3945941379}.{} (letters from June 1810 to June 1854)}
\item
  {Clemens Schaefer, ed. (1927). Briefwechsel zwischen Carl Friedrich Gauss und Christian Ludwig Gerling (in German). Berlin: Otto Elsner.{} (letters from June 1810 to June 1854)}
\item
  {\href{https://en.wikipedia.org/wiki/Karl_Christian_Bruhns}{Karl Christian Bruhns}, ed. (1877). \href{https://archive.org/details/briefezwischena00gausgoog/page/n5/mode/2up}{\emph{Briefe zwischen A. v. Humboldt und Gauss}} (in German). Leipzig: Wilhelm Engelmann.{} (letters from July 1807 to December 1854; added letters of other correspondents)}
\item
  {\href{https://en.wikipedia.org/wiki/Karin_Reich}{Reich, Karin}; Roussanova, Elena (2018). \emph{\href{https://en.wikipedia.org/wiki/Karl_Kreil}{Karl Kreil} und der Erdmagnetismus. Seine Korrespondenz mit Carl Friedrich Gauß im historischen Kontext}. Veröffentlichungen der Kommission für Geschichte der Naturwissenschaften, Mathematik und Medizin, No. 68 (in German). Vienna: Verlag der Österreichischen Akademie der Wissenschaften.{} (letters from 1835 to 1843)}
\item
  {Gerardy, Theo, ed. (1959). \emph{Briefwechsel zwischen Carl Friedrich Gauß und \href{https://en.wikipedia.org/wiki/Karl_Ludwig_von_Lecoq}{Carl Ludwig von Lecoq}}. Abhandlungen der Akademie der Wissenschaften in Göttingen, Mathematisch-Physikalische Klasse, No. 4 (in German). Göttingen: Vandenhoeck \& Ruprecht. pp.{~}{37--}63.{} (letters from February 1799 to September 1800)}
\item
  {Forbes, Eric G. (1971). \href{https://www.tandfonline.com/doi/abs/10.1080/00033797100203767}{"The Correspondence between Carl Friedrich Gauss and the Rev. Nevil Maskelyne (1802--05)"}. \emph{Annals of Science}. \textbf{27} (3): {213--}237. \href{https://en.wikipedia.org/wiki/Doi_(identifier)}{doi}:\href{https://doi.org/10.1080\%2F00033797100203767}{10.1080/00033797100203767}.{}}
\item
  {\href{https://en.wikipedia.org/wiki/Clifford_Cunningham}{Cunningham, Clifford} (2004). \href{https://www.tandfonline.com/doi/abs/10.1080/00033790310001660164}{"Discovery of the Missing Correspondence between Carl Friedrich Gauss and the Rev. Nevil Maskelyne (1802--05)"}. \emph{Annals of Science}. \textbf{61} (4): {469--}481. \href{https://en.wikipedia.org/wiki/Doi_(identifier)}{doi}:\href{https://doi.org/10.1080\%2F00033790310001660164}{10.1080/00033790310001660164}.{}}
\item
  {Carl Schilling, ed. (1900). \href{https://archive.org/details/p1wilhelmolberss02olbeuoft/page/n5/mode/2up}{\emph{Briefwechsel zwischen Olbers und Gauss: Erste Abtheilung}}. \href{https://en.wikipedia.org/wiki/Wilhelm_Olbers}{Wilhelm Olbers}. Sein Leben und seine Werke. Zweiter Band (in German). Berlin: Julius Springer.{} (letters from January 1802 to October 1819)}
\item
  {Carl Schilling, ed. (1909). \href{https://archive.org/details/p2wilhelmolberss02olbeuoft/page/n7/mode/2up}{\emph{Briefwechsel zwischen Olbers und Gauss: Zweite Abtheilung}}. Wilhelm Olbers. Sein Leben und seine Werke. Zweiter Band (in German). Berlin: Julius Springer.{} (letters from January 1820 to May 1839; added letters of other correspondents)}
\item
  {\href{https://en.wikipedia.org/wiki/Christian_August_Friedrich_Peters}{Christian August Friedrich Peters}, ed. (1860--1865). \emph{Briefwechsel zwischen C. F. Gauss und \href{https://en.wikipedia.org/wiki/H._C._Schumacher}{H. C. Schumacher}} (in German). Altona: Gustav Esch.{}}

  \begin{itemize}
  \tightlist
  \item
    {\href{https://books.google.com/books?id=aIRtAAAAMAAJ}{Volumes 1+2} (letters from April 1808 to March 1836)}
  \item
    {\href{https://books.google.com/books?id=NDcDAAAAQAAJ}{Volumes 3+4} (letters from March 1836 to April 1845)}
  \item
    {\href{https://books.google.com/books?id=3jEDAAAAQAAJ}{Volumes 5+6} (letters from April 1845 to November 1850)}
  \end{itemize}
\item
  {Poser, Hans, ed. (1987). \emph{Briefwechsel zwischen Carl Friedrich Gauß und \href{https://en.wikipedia.org/wiki/Eberhard_August_Wilhelm_Zimmermann}{Eberhard August Zimmermann}}. Abhandlungen der Akademie der Wissenschaften in Göttingen, Mathematisch-Physikalische Klasse, Folge 3, No. 39 (in German). Göttingen: Vandenhoeck \& Ruprecht. \href{https://en.wikipedia.org/wiki/ISBN_(identifier)}{ISBN}{~}\href{https://en.wikipedia.org/wiki/Special:BookSources/978-3525821169}{978-3525821169}.{} (letters from 1795 to 1815)}
\end{itemize}

The \href{https://en.wikipedia.org/wiki/Göttingen_Academy_of_Sciences_and_Humanities}{Göttingen Academy of Sciences and Humanities} provides a complete collection of the known letters from and to Carl Friedrich Gauss that is accessible online. The literary estate is kept and provided by the \href{https://en.wikipedia.org/wiki/Göttingen_State_and_University_Library}{Göttingen State and University Library}. Written materials from Carl Friedrich Gauss and family members can also be found in the municipal archive of Brunswick.





\begin{itemize}
\tightlist
\item
  \href{https://en.wikipedia.org/wiki/Uta_Merzbach}{Merzbach, Uta C.}; \href{https://en.wikipedia.org/wiki/Carl_Benjamin_Boyer}{Boyer, Carl B.} (2011). \href{https://books.google.com/books?id=bR9HAAAAQBAJ}{\emph{A History of Mathematics}} (3rd{~}ed.). New Jersey: John Wiley \& Sons. \href{https://en.wikipedia.org/wiki/ISBN_(identifier)}{ISBN}{~}\href{https://en.wikipedia.org/wiki/Special:BookSources/978-0470630563}{978-0470630563}.{}
\item
  \href{https://en.wikipedia.org/wiki/Tord_Hall}{Hall, Tord} (1970). \href{https://archive.org/details/carlfriedrichgau00tord}{\emph{Carl Friedrich Gauss: A Biography}}. Cambridge, MA: \href{https://en.wikipedia.org/wiki/MIT_Press}{MIT Press}. \href{https://en.wikipedia.org/wiki/ISBN_(identifier)}{ISBN}{~}\href{https://en.wikipedia.org/wiki/Special:BookSources/978-0-262-08040-8}{978-0-262-08040-8}. \href{https://en.wikipedia.org/wiki/OCLC_(identifier)}{OCLC}{~}\href{https://search.worldcat.org/oclc/185662235}{185662235}.{}
\item
  \href{https://en.wikipedia.org/wiki/Paul_J._Nahin}{Nahin, Paul J.} (2010). \href{https://books.google.com/books?id=OPyPwaElDvUC&pg=PA82}{\emph{An Imaginary Tale: The Story of √-1}}. Princeton University Press. \href{https://en.wikipedia.org/wiki/ISBN_(identifier)}{ISBN}{~}\href{https://en.wikipedia.org/wiki/Special:BookSources/9_78-1-4008-3389-4}{9 78-1-4008-3389-4}.{}
\item
  Simmons, J. (1996). \emph{The Giant Book of Scientists: The 100 Greatest Minds of All Time}. Sydney: The Book Company.{}
\item
  Tent, Margaret (2006). \emph{The Prince of Mathematics: Carl Friedrich Gauss}. A. K. Peters. \href{https://en.wikipedia.org/wiki/ISBN_(identifier)}{ISBN}{~}\href{https://en.wikipedia.org/wiki/Special:BookSources/978-1-56881-455-1}{978-1-56881-455-1}.{}
\end{itemize}


\subsubsection{Fictional}\label{Fictional}

\begin{itemize}
\tightlist
\item
  \href{https://en.wikipedia.org/wiki/Daniel_Kehlmann}{Kehlmann, Daniel} (2005). \href{https://en.wikipedia.org/wiki/Measuring_the_World}{\emph{Die Vermessung der Welt}} (in German). Rowohlt. \href{https://en.wikipedia.org/wiki/ISBN_(identifier)}{ISBN}{~}\href{https://en.wikipedia.org/wiki/Special:BookSources/978-3-498-03528-0}{978-3-498-03528-0}. \href{https://en.wikipedia.org/wiki/OCLC_(identifier)}{OCLC}{~}\href{https://search.worldcat.org/oclc/144590801}{144590801}.{}

  \begin{itemize}
  \tightlist
  \item
    Kehlmann, Daniel (2006). \emph{\href{https://en.wikipedia.org/wiki/Measuring_the_World}{Measuring the World}}. Translated by \href{https://en.wikipedia.org/wiki/Carol_Brown_Janeway}{Janeway, Carol Brown}.{}
  \end{itemize}
\end{itemize}



{\href{https://en.wikipedia.org/wiki/File:Commons-logo.svg}{
\includegraphics[width=0.3125in,height=0.41667in]{40px-Commons-logo.svg.png}}}

Wikimedia Commons has media related to \href{https://commons.wikimedia.org/wiki/Johann\%20Carl\%20Friedrich\%20Gauß}{{Johann Carl Friedrich Gauß}}.


{\href{https://en.wikipedia.org/wiki/File:Wikisource-logo.svg}{
\includegraphics[width=0.39583in,height=0.41667in]{40px-Wikisource-logo.svg.png}}}

English \href{https://en.wikipedia.org/wiki/Wikisource}{Wikisource} has original text related to this article:

\textbf{\href{https://en.wikisource.org/wiki/en:Author:Carl\%20Friedrich\%20Gauss}{Carl Friedrich Gauss}}


{{
\includegraphics[width=0.35417in,height=0.41667in]{40px-Wikiquote-logo.svg.png}}}

Wikiquote has quotations related to \emph{\textbf{\href{https://en.wikiquote.org/wiki/Special:Search/Carl\%20Friedrich\%20Gauss}{Carl Friedrich Gauss}}}.


\begin{itemize}
\tightlist
\item
  Works by or about Carl Friedrich Gauss at the \href{https://en.wikipedia.org/wiki/Internet_Archive}{Internet Archive}
\item
  Publications of C. F. Gauss in \href{https://en.wikipedia.org/wiki/Astrophysics_Data_System}{Astrophysics Data System}
\item
  \href{http://adsabs.harvard.edu//full/seri/MNRAS/0016//0000080.000.html}{"Obituary"}. \emph{Monthly Notices of the Royal Astronomical Society}. \textbf{16} (80). 1856{. Retrieved {8 April} 2024}.{}
\item
  \href{https://www.bbc.co.uk/programmes/b00ss0lf}{"Carl Friedrich Gauss"}. \emph{BBC 4}. A Brief History of Mathematics{. Retrieved {8 April} 2024}.{}
\item
  \href{http://american_almanac.tripod.com/gauss.htm}{"Mind Over Mathematics: How Gauss Determined The Date of His Birth"}. \emph{american\_almanac.tripod.com}{. Retrieved {8 April} 2024}.{}
\item
  \href{https://mathoverflow.net/questions/370190/how-did-gauss-characterize-the-metrical-relations-in-the-uniform-4-4-4-tiling}{"On Gauss\textquotesingle{} tesselation of the unit disk"}{. Retrieved {8 April} 2024}.{}
\item
  \href{https://www.uni-goettingen.de/en/62596.html}{"Carl Friedrich Gauß"}. Göttingen University{. Retrieved {8 April} 2024}.{}
\item
  " Descendants of Gauss ". \emph{www.coloradocollege.edu}. Archived from \href{http://www.coloradocollege.edu/library/SpecialCollections/CenturyChest/transcription24.html}{the original} on 28 May 2010{. Retrieved {4 June} 2024}.{}
\item
  \href{http://www.gausschildren.org/cfgauss/cfgauss.htm}{"Gauss"}. \emph{www.gausschildren.org}{. Retrieved {8 April} 2024}.{}
\item
  \href{https://denkmalatlas.niedersachsen.de/viewer/themen/Gauss-Steine/}{Carl Friedrich Gauss -- Spuren seines Lebens} (Places used as points for triangulation)
\end{itemize}









{}
